\documentclass[12pt, letterpaper, oneside]{book}  
\usepackage[spanish, es-tabla, es-lcroman]{babel}
\usepackage[top=2.5cm, bottom=2.5cm, left=3cm, right=2cm]{geometry}
\usepackage[utf8]{inputenc} 
\usepackage{amsmath}
\usepackage{amssymb}
\usepackage{dsfont}
\usepackage{graphicx}
\usepackage{color}
\usepackage{etoolbox}
\usepackage{titlesec}
\usepackage{blindtext}
\usepackage{multirow}
\usepackage{pdfpages}
\usepackage{fancyhdr}
\usepackage{hyperref}
\usepackage{pgfplots}
\usepackage{rotating}
\usepackage{multirow}
\usepackage{multicol}
\usepackage{paralist}
\usepackage{longtable}
\usepackage{subcaption}
\usepackage{listings}
\usepackage{epigraph}
\usepackage{setspace}
\onehalfspacing
\decimalpoint
\usepackage[comma, square, sort, numbers]{natbib}
\usepackage[font=small,labelfont=bf]{caption}
\renewcommand{\sin}{\sen}
\newcommand*{\unnref}[1]{\hyperref[{#1}]{\nameref*{#1}}}

\definecolor{mygreen}{rgb}{0,0.6,0}
\definecolor{mygray}{rgb}{0.5,0.5,0.5}
\definecolor{mymauve}{rgb}{0.58,0,0.82}

\lstset{literate=
	{á}{{\'a}}1 {é}{{\'e}}1 {í}{{\'i}}1 {ó}{{\'o}}1 {ú}{{\'u}}1
	{Á}{{\'A}}1 {É}{{\'E}}1 {Í}{{\'I}}1 {Ó}{{\'O}}1 {Ú}{{\'U}}1
	{à}{{\`a}}1 {è}{{\`e}}1 {ì}{{\`i}}1 {ò}{{\`o}}1 {ù}{{\`u}}1
	{À}{{\`A}}1 {È}{{\'E}}1 {Ì}{{\`I}}1 {Ò}{{\`O}}1 {Ù}{{\`U}}1
	{ä}{{\"a}}1 {ë}{{\"e}}1 {ï}{{\"i}}1 {ö}{{\"o}}1 {ü}{{\"u}}1
	{Ä}{{\"A}}1 {Ë}{{\"E}}1 {Ï}{{\"I}}1 {Ö}{{\"O}}1 {Ü}{{\"U}}1
	{â}{{\^a}}1 {ê}{{\^e}}1 {î}{{\^i}}1 {ô}{{\^o}}1 {û}{{\^u}}1
	{Â}{{\^A}}1 {Ê}{{\^E}}1 {Î}{{\^I}}1 {Ô}{{\^O}}1 {Û}{{\^U}}1
	{Ã}{{\~A}}1 {ã}{{\~a}}1 {Õ}{{\~O}}1 {õ}{{\~o}}1
	{œ}{{\oe}}1 {Œ}{{\OE}}1 {æ}{{\ae}}1 {Æ}{{\AE}}1 {ß}{{\ss}}1
	{ű}{{\H{u}}}1 {Ű}{{\H{U}}}1 {ő}{{\H{o}}}1 {Ő}{{\H{O}}}1
	{ç}{{\c c}}1 {Ç}{{\c C}}1 {ø}{{\o}}1 {å}{{\r a}}1 {Å}{{\r A}}1
	{€}{{\euro}}1 {£}{{\pounds}}1 {«}{{\guillemotleft}}1
	{»}{{\guillemotright}}1 {ñ}{{\~n}}1 {Ñ}{{\~N}}1 {¿}{{?`}}1
}
\lstset{ 
	backgroundcolor=\color{white},
	basicstyle=\footnotesize,
	breakatwhitespace=false,
	breaklines=true,
	captionpos=b,
	commentstyle=\color{mygreen},
	deletekeywords={...},
	escapeinside={\%*}{*)},
	extendedchars=true,
	firstnumber=1,
	frame=single,
	keepspaces=true,
	keywordstyle=\color{blue},
	language=Python,
	morekeywords={*,...},
	numbers=left,
	numbersep=5pt,
	numberstyle=\tiny\color{mygray},
	rulecolor=\color{black},
	showspaces=false,
	showstringspaces=false,
	showtabs=false,
	stepnumber=1,
	stringstyle=\color{mymauve},
	tabsize=2,
	title=\lstname
}
\lstset{ 
	backgroundcolor=\color{white},
	basicstyle=\footnotesize,
	breakatwhitespace=false,
	breaklines=true,
	captionpos=b,
	commentstyle=\color{mygreen},
	deletekeywords={...},
	escapeinside={\%*}{*)},
	extendedchars=true,
	firstnumber=1,
	frame=single,
	keepspaces=true,
	keywordstyle=\color{blue},
	language=sh,
	morekeywords={*,...},
	numbers=left,
	numbersep=5pt,
	numberstyle=\tiny\color{mygray},
	rulecolor=\color{black},
	showspaces=false,
	showstringspaces=false,
	showtabs=false,
	stepnumber=1,
	stringstyle=\color{mymauve},
	tabsize=2,
}
\lstset{ 
	backgroundcolor=\color{white},
	basicstyle=\footnotesize,
	breakatwhitespace=false,
	breaklines=true,
	captionpos=b,
	commentstyle=\color{mygreen},
	deletekeywords={...},
	escapeinside={\%*}{*)},
	extendedchars=true,
	firstnumber=1,
	frame=single,
	keepspaces=true,
	keywordstyle=\color{blue},
	language=c++,
	morekeywords={*,...},
	numbers=left,
	numbersep=5pt,
	numberstyle=\tiny\color{mygray},
	rulecolor=\color{black},
	showspaces=false,
	showstringspaces=false,
	showtabs=false,
	stepnumber=1,
	stringstyle=\color{mymauve},
	tabsize=2,
}
\setlength\epigraphwidth{.5\textwidth}
\setlength\epigraphrule{0pt}

\hypersetup{colorlinks=true, linkcolor=black, citecolor=black, urlcolor=black}
\renewcommand{\labelitemi}{$\blacktriangleright$}
\renewcommand{\labelitemii}{$\vartriangleright$}
\definecolor{gray75}{gray}{0.75}
\titlespacing*{\chapter}{0pt}{-15pt}{20pt}	
\newcommand{\hsp}{\hspace{20pt}}
\titleformat{\chapter}[hang]{\Huge\bfseries}{\thechapter{.}\hsp}{0pt}{\Huge\bfseries}
\fancypagestyle{plain}
{
	\fancyhf{}
	\rhead{\includegraphics[width=0.75cm]{./Figuras/UPIIZ}}
	\chead{IPN-UPIIZ}
	\lhead{\includegraphics[width=0.75cm]{./Figuras/IPN}}
	\rfoot{\thepage}
	\renewcommand{\headrulewidth}{0.4pt}
	\renewcommand{\footrulewidth}{0pt}
}
\pagestyle{fancy}
	\fancyhf{}
	\rhead{\includegraphics[width=0.75cm]{./Figuras/UPIIZ}}
	\chead{IPN-UPIIZ}
	\lhead{\includegraphics[width=0.75cm]{./Figuras/IPN}}
	\rfoot{\thepage}
	\renewcommand{\headrulewidth}{0.4pt}
	\renewcommand{\footrulewidth}{0pt}

\begin{document}
	\thispagestyle{empty}
	\includepdf[pages=-]{./Portada/Portada.pdf}
	\thispagestyle{empty}
	\includepdf[pages=-]{./Portada2/Portada2.pdf}
	\thispagestyle{empty}
	\begin{center}
		{\bf Autorización de uso de obra\\
			INSTITUTO POLITÉCNICO NACIONAL\\
			Presente}\\
	\end{center}
	
	Bajo protesta de decir la verdad el que suscribe, {\bf Claudio Iván Esparza Castañeda}, estudiante del programa de {\bf Ingeniería Mecatrónica}, con número de boleta {\bf 2016670022}, adscrito a la {\bf Unidad Profesional Interdisciplinaria de Ingeniería Campus Zacatecas}; manifiesto ser autor y titular de los derechos morales y patrimoniales de la obra titulada {\bf ``Implementación de algoritmos de conducción autónoma en vehículos a escala 1:10''}, en adelante el ``Trabajo de Titulación'' y de la cual se adjunta copia, por lo que por medio del presente y con fundamento en el artículo 27 fracción II, inciso b) de la Ley Federal del Derecho de autor, otorgamos al Instituto Politécnico Nacional, en adelante IPN, autorización no exclusiva para comunicar y exhibir públicamente total o parcialmente en medios digitales el ``Trabajo de Titulación'' por un periodo de tiempo indefinido contado a partir de la fecha de la presente autorización, dicho periodo se renovará automáticamente en caso de no dar aviso expreso a el IPN de su terminación.
	\par En virtud de lo anterior, el IPN deberá reconocer en todo momento mi autoría en el ``Trabajo de Titulación''.
	\par Adicionalmente, y en mi calidad de autor y titular de los derechos morales y patrimoniales de el ``Trabajo de Titulación'', manifiesto que la misma es original y que la presente autorización no contraviene ninguna otorgada por los suscritos respecto del ``Trabajo de Titulación'', por lo que deslindo de toda responsabilidad a el IPN en caso de que el contenido de el ``Trabajo de Titulación'' o la autorización concedida afecte o viole derechos autorales, industriales, secretos industriales, convenios o contratos de confidencialidad o en general cualquier derecho de propiedad intelectual de terceros y asumo las consecuencias legales y económicas de cualquier demanda o reclamación que pueda derivarse del caso.
	\begin{flushright}
		Zacatecas, Zac., a 14 de Mayo del 2021
	\end{flushright}
	\begin{center}
		Atentamente
		\vspace*{1cm}\\
		\rule{7cm}{1pt}\\
		Claudio Iván Esparza Castañeda
	\end{center}
	\newpage
	\thispagestyle{empty}
	\epigraph{'má, te dedico este Trabajo Terminal y un verso: {\it<<Te lo agradezco, madre: hay que seguir levantándose temprano para esperar diariamente la vida>>.}}{Estés donde estés, siempre estarás aquí.}
	\newpage
	\thispagestyle{empty}
	Estoy agradecido por toda la eternidad con mis padres Juan y María, quienes me han apoyado de manera incondicional durante toda mi etapa de estudiante y en general en mi vida. Quiero extender este agradecimiento a mis hermanos Karen y Abraham, tíos y tías, primas y primos, abuelos y abuelas.
	\par Levanto la voz en un {\it Huelum} a mi {\it alma mater} el Instituto Politécnico Nacional y otro a la Unidad Profesional Interdisciplinaria de Ingeniería Campus Zacatecas, por brindarme los conocimientos que permitieron este proyecto, mismos estudios que después serán mi bastión durante toda mi vida profesional, por formarme cada día para así poder poner la técnica al servicio de la patria. Un {\it Huelum} a mis maestros, por ser los proveedores directos de esos conocimientos. Un {\it Huelum} a mis asesores, Flabio y Carlos Daniel, por resolver mis infinitas dudas y orientarme a lo largo del desarrollo de este proyecto. Gracias totales.
	\par Un último {\it Huelum} a mis compañeros a quienes ahora considero como mis amigos al transitar a la par el camino durante estos años: Ronaldo, Miguel, Eduardo, Julio, Ulises, Daniel; el mayor de los éxitos, colegas. A Alejandra, que desde la preparatoria me ha impulsado a seguir creciendo, un abrazo y muchas gracias. A los colegas divulgadores del {\it ZigZag}: Iván, Blanca, Corina, Juan Carlos, Armando, Alondra, Jimena, Pepe, Diana, Jonathan, Osvaldo, Victoria, Mayra, etc.; que siempre han propiciado mi aprendizaje a través de sus distintos puntos de vista. A mis amigos del {\it CEBAARE}: Gandhi, Joaquín, Brandon, Isaac, Alan, Ramcés, Gerardo, Misael y Gechén (el honorario); sí vamos a poder y estamos pudiendo. A todos ellos, gracias.
	\par ¡Politécnico, gloria!
	\chapter*{Resumen}
	\label{cap:res}
	\addcontentsline{toc}{chapter}{Resumen}
	\pagenumbering{roman}
	En este documento se reúne el proceso de desarrollo de un prototipo de vehículo autónomo a escala 1:10. El prototipo se desplaza sobre una pista de pruebas que simula un entorno urbano. La plataforma robótica está compuesta por un chasís comercial que lleva sobre sí un sistema de visión artificial basado en un sensor tipo {\it Kinect}, mismo que a su vez se apoya en una mini-computadora a bordo para procesar la información visual recibida, permitiéndole reconocer de esta manera las características de las líneas del carril sobre el que se desplaza el vehículo, además de realizar el reconocimiento de las señales de Alto, vuelta a la Izquierda, vuelta a la Derecha y la identificación del color en las luces de un semáforo; permitiéndole controlar al vehículo permaneciendo dentro del carril y actuar de acuerdo con las señales de tránsito detectadas. Por las características de la computadora, no es posible enviar las señales de PWM a los motores de manera directa con la misma, por lo que para realizar esta acción la computadora se apoya en un sistema basado en microcontrolador.
	\section*{Palabras clave}
	\label{sec:palcla}
	\addcontentsline{toc}{section}{Palabras clave}
	Robótica Móvil, Navegación Autónoma, Visión Artificial, Control Visual.
	\chapter*{Abstract}
	\label{cap:abs}
	\addcontentsline{toc}{chapter}{Abstract}
	This document is a compilation of the development process for a 1:10 scaled autonomous car. The prototype moves on a test track that simulates a urban environment. The robotic platform is made up by a comercial chassis which carries on itself a {\it Kinect} type based computer vision system, the images taken by the RGB-D sensor are processed by an on board mini-computer, recognizing by this way the Stop, turn Left, turn Right signs and the light color identification in a traffic light; allowing it control the vehicle keeping the lane and acting according to the road signs detected. By the computer features, isn't possible send the PWM signals directly by itself, this is the reason why a microcontroller-based system helps the computer in this action.
	\section*{Keywords}
	\label{sec:kw}
	\addcontentsline{toc}{section}{Keywords}
	Mobile Robotics, Autonomous Navigation, Computer Vision, Visual-Feedback Control.
	\renewcommand{\listfigurename}{Lista de Figuras}
	\renewcommand{\listtablename}{Lista de Tablas}
	\renewcommand{\contentsname}{Índice General}
	\renewcommand*{\lstlistlistingname}{Lista de Códigos}
	\tableofcontents
	\cleardoublepage
	\addcontentsline{toc}{chapter}{Lista de Figuras}
	\listoffigures
	\cleardoublepage
	\addcontentsline{toc}{chapter}{Lista de Tablas} 
	\listoftables
	\cleardoublepage
	\addcontentsline{toc}{chapter}{Lista de Códigos} 
	\lstlistoflistings
	\chapter*{Introducción}
	\label{cap:intro}
	\addcontentsline{toc}{chapter}{Introducción}
	\pagenumbering{arabic}
	Uno de los problemas principales que existe dentro del campo de la robótica móvil es la navegación autónoma. Los sistemas capaces de efectuar una navegación autónoma robusta y eficiente tiene varias aplicaciones tanto industriales como domésticas, siendo el desarrollo de vehículos totalmente autónomos una de las aplicaciones con mayores avances y expectativas. Los principales progresos realizados en esta área de investigación se deben a las universidades y los centros de desarrollo de las compañías automotrices, en gran parte ha sido posible gracias a la integración de una gran variedad de sensores; entre estos, los sistemas de visión artificial han cobrado protagonismo por su capacidad de proveer una considerable cantidad de información a un costo relativamente bajo; gracias a los mismos es que se han desarrollado técnicas de control con realimentación visual, mapeo y reconstrucción de espacios y odometría visual.
	\par Los vehículos autónomos se encontraban relegados a la ciencia ficción. Hasta hace poco comenzaron a salir de las novelas y películas de ciencia ficción para pasar al escenario de la realidad. Enumerando de manera rápida algunos de los principales avances en este campo remiten al {\it DARPA Grand Challenge} en marzo de 2004, donde vehículos desarrollados por un gremio conformado por centros de investigación de universidades, fabricantes de autos y hasta aficionados se reunieron para desarrollar un vehículo que sin conductor pudiera completar una ruta de 150 millas sin intervención humana; sin embargo ninguno de los vehículos participantes completó el desafío; la historia cambió para la segunda edición de este reto en el mes de octubre del 2005, donde cinco de los vehículos competidores completaron de manera exitosa 132 millas del recorrido.
	\par Durante los años que han transcurrido desde estos acontecimientos, la cantidad de estudios relacionados con la conducción autónoma ascendió, teniendo como panorama actual a una gran cantidad de estos vehículos que están utilizando a las carreteras y a las ciudades como pista de pruebas. En varias ciudades del mundo los vehículos autónomos circulan a la par que se mejoran las tecnologías que involucran. Parte de los más enérgicos competidores son las compañías automotrices como Tesla, Ford o Chevrolet. Impulsados en gran medida por estas compañías, existen competencias similares al DARPA Grand Challenge enfocadas a promover el desarrollo de vehículos autónomos a pequeña escala, mismos que pretenden emular las prestaciones de los vehículos de tamaño real; entre los beneficios de utilizar sistemas a pequeña escala por sobre los sistemas a tamaño completo, está de más mencionar el ahorro de los costos relacionados al propio vehículo, los sistemas sensoriales y de cómputo a bordo, y también al escenario de pruebas que es necesario. 
	\par El vehículo autónomo del que es objeto este Trabajo Terminal fue motivado en primer lugar por el {\it AutoNOMOS Model}, desarrollado por el Dr. Raúl Rojas y su equipo de la {\it Freie Universität Berlin}, o los construidos por el {\it MIT} o la {\it University of Pennsylvania}, entre tantos otros que se expolorarán en el el \autoref{cap:eda}. Este trabajo aborda la navegación autónoma de un vehículo a escala 1:10 de manera reactiva, lo que significa que el carro pueda permanecer en el carril por el que se desplaza y cumplir con la identificación de los señalamientos de tránsito gracias a la realimentación visual provista por el sistema de visión presente en el robot, que a su vez se apoya en la computadora a bordo para los cálculos y el procesamiento de la información recibida.
	\par Antes de describir la estructura general de este documento, se aclara por adelantado que a este robot también se denominará de de diferentes maneras a lo largo del documento, como RMR\footnote{Robot Móvil con Ruedas}, plataforma, vehículo, carro o coche.
	\par Los contenidos que proceden a esta Introducción se dan de la siguiente manera: En dos capítulos no numerados permanecen inmutables desde el Protocolo de Investigación los \unnref{cap:obj} y la \unnref{cap:jus}. En el \autoref{cap:mt} se destaca la limitación teórica que da sustento al trabajo, tales como la robótica móvil, las perspectivas de la navegación autónoma, la visión por computadora y las técnicas de control aplicadas a la navegación autónoma. El \autoref{cap:eda} trata acerca de los estudios que anteceden en tiempo y técnica a este Trabajo Terminal, y que sirven como base de lo aquí desarrollado. A lo largo del \autoref{cap:pdp} se muestra una perspectiva de lo que en primera instancia se ha planteado  lograr con este trabajo, explicando las tecnologías que a priori se pensó emplear y se da una primera propuesta de la pista de pruebas. En el \autoref{cap:des} se da una extensa descripción de la metodología y los trabajos realizados para lograr concluir de manera satisfactoria el trabajo, tomando como elementos al modelo matemático del sistema, las propuestas de control, la selección de los componentes, el ensamblaje de la plataforma y también la construcción de la pista de pruebas, quedando como elemento final de este capítulo la descripción e implementación de los algoritmos de navegación autónoma que forman parte del proyecto. En el \autoref{cap:resul} se realiza un análisis y discusión de los resultados obtenidos, con lo cual de manera crítica se permitirá desembocar en el \autoref{cap:con}, donde se dan conclusiones generales acerca de los resultados y se plantean nuevos retos como posible trabajo futuro. Como adendum se incluye aquel material adicional que fue de ayuda y que por su valor documental es importante su incorporación para futuras referencias bibliográficas.
	\chapter*{Objetivos}
	\label{cap:obj}
	\addcontentsline{toc}{chapter}{Objetivos}
	\section*{Objetivo general}
	\label{sec:objg}
	\addcontentsline{toc}{section}{Objetivo general}
	Implementar algoritmos de visión artificial para la identificación de características del entorno, que serán usadas como realimentación visual en un esquema de control para la navegación de un vehículo autónomo a escala en un escenario urbano simulado.
	\section*{Objetivos específicos}
	\label{sec:obje}
	\addcontentsline{toc}{section}{Objetivos específicos}
	\begin{itemize}
		\item Escoger una plataforma robótica de propulsión eléctrica para validar la navegación autónoma del proyecto.
		\item Seleccionar los algoritmos de visión artificial adecuados para el proyecto.
		\item Realizar el análisis cinemático del robot móvil para la obtención de parámetros que luego se emplearán en el control del vehículo.
		\item Diseñar y programar los algoritmos de control visual para su implementación en el robot móvil.
		\item Usar técnicas de control para regular la velocidad y dirección del vehículo autónomo a escala.
		\item Emplear la realimentación visual para alcanzar un nivel SAE 3 de automatización\cite{saeTaxonomyDefinitiosTerms2014} en el vehículo a escala.
		\item Construir un escenario de pruebas para la validación de la plataforma a través de:
		\begin{itemize}
			\item Identificar y permanecer en el carril por el que se desplaza.
			\item Reconocer las señales  de alto, vuelta a la izquierda y vuelta a la derecha en el escenario de pruebas.
			\item Examinar el color de la luz en un semáforo y reaccionar de manera correcta ante ello.	
		\end{itemize}
	\end{itemize}
	\chapter*{Justificación}
	\label{cap:jus}
	\addcontentsline{toc}{chapter}{Justificación}
	La necesidad principal de la autonomía en la robótica móvil, así como en la industria automotriz, se da cuando un operador humano no puede maniobrar el vehículo de manera directa o teleoperada. Para la conducción autónoma, se pretende reducir el error humano asociado a la conducción.
	\par En la actualidad, existen cientos de vehículos en fase de pruebas por compañías del ramo automotriz y tecnológico \cite{Autopilot, LookingFurther, WaymoWaymo}. Estas pruebas requieren, en su primera etapa, una ciudad construida con el propósito de afinar las tecnologías antes de llevar el vehículo a las calles de una ciudad real. Los costos del vehículo per se, añadiendo todos los sistemas que lo ayudan a guiarse (cámara monocular o estereoscópica, sensores para la detección de obstáculos, computadora a bordo, etc.), se ven incrementados por el costo que representa construir un escenario, con las características de una ciudad, que funcione con el único propósito de proceder con la realización de las pruebas antes de llevar el vehículo a una urbe, con todo lo que esto representa. El surgimiento de competencias para la validación de los sistemas automotrices inteligentes \cite{GrandChallengeAutonomous}, aunado al interés de las gigantes automotrices en las competiciones de vehículos autónomos a escala, es un indicativo de que es necesario optimizar los algoritmos empleados, teniendo como opción realizarlos antes en un entorno simulado y después en la ciudad. Es por esta situación que se tiene el propósito de cambiar el vehículo de tamaño real por un vehículo a pequeña escala, que pueda emplear los sistemas antes mencionados, en un esquema de control con realimentación visual, sobre un escenario de pruebas.
	\chapter{Marco teórico}
	\label{cap:mt}
	En este Capítulo se abordan los conceptos básicos en los que se sustenta este trabajo, se busca mostrar de esta manera la postura conceptual de este objeto de estudio. En primer lugar, en la Sección \ref{sec:rm} se exploran los formalismos relacionados a la robótica móvil, además de los sistemas de medición como parte de la percepción de los estados internos y externos del robot aplicado al mismo. Para la Sección \ref{sec:na} se presentan un par de avances históricos importantes de la navegación autónoma dentro de la robótica móvil y las perspectivas actuales de la misma junto con los propios esquemas que buscan cuantificar la autonomía en los vehículos. La Sección \ref{sec:vi} se enfoca en la visión artificial, y en particular se exploran al final los sistemas de visión tridimensional, haciendo énfasis en los sensores RGB-D, tomando como punto de partida su uso dentro de la robótica móvil. Para finalizar el capítulo, en la Sección \ref{sec:con} se presentan las estrategias de control más usadas para este tipo de vehículos, se plantean los principales enfoques y se encuadra al {\it visual servoing} junto con sus características primarias.
\section{Robótica móvil}
\label{sec:rm}
A pesar de que existen un gran número de registros históricos acerca de autómatas que fueron elaborados por personajes que van desde Herón de Alejandría hasta Leonardo Da Vinci, la robótica es una rama del conocimiento joven en realidad, pues no se consolidó como ciencia hasta el siglo XX. En primera instancia la palabra robot se acuñó en la obra {\it Rossum's Universal Robots} de Carel Capek, la palabra proviene del idioma checo y en ese idioma se suele usar de manera coloquial para designar a un sirviente; en la cultura popular surgieron lo robots como personajes que podían realizar las tareas que a los humanos llegan a resultar de gran dificultad, y con la aparición de las tres leyes de la robótica de Isaac Asimov en {\it I, Robot}, se establecían las pautas a las que los fabricantes debían someter a los robots de la vida real.
\par La robótica se divide en las ramas de robótica de manipuladores y robótica móvil, siendo esta última la de interés en este trabajo. Desde el punto de vista más general, se define al robot móvil como aquel que es capaz de cambiar su posición sin estar sujeto a un solo punto. Las aplicaciones de los robots móviles son en realidad amplias, pues pueden ser usados en tareas de exploración, búsqueda y rescate, automatización de procesos, vigilancia, investigación científica y agricultura, por sólo mencionar algunas. Los robots, antes de alcanzar su grado de desarrollo actual, atravesaron un largo proceso de evolución y desarrollo. De ser mecanismos teleoperados, han llegado a adquirir rasgos que les brindan autonomía a través de acciones que les permiten medir características del entorno, calculando y actuando a partir de ello.
\subsection{Locomoción}
\label{ssec:loc}
Existen diferentes posibilidades de lograr el desplazamiento de un robot móvil, incluyendo (pero no limitando): Rodar, caminar, saltar, trepar, nadar, volar, etc. De acuerdo al medio en el que se desplaza se pueden clasificar en: Acuático, terrestre, aéreo e híbrido (combinando al menos dos de los medios).
\begin{table}[htbp!]
	\caption{Tipos de ruedas empleados en robótica móvil.}
	\label{tab:ruedas}
	\begin{center}
		\begin{tabular}{|p{3cm}|p{7cm}|p{4cm}|}
			\hline
			{\bf Tipo de rueda} & {\bf Grados de libertad} & {\bf Ejemplo}\\
			\hline
			Estándar \includegraphics[width=2.5cm]{./Figuras/REstan.png} & Dos:
			\begin{compactitem}
				\item Rotación alrededor de su eje.
				\item Rotación alrededor de su punto de contacto con el suelo.
			\end{compactitem} & Rueda de carretilla\\
			\hline
			Castor \includegraphics[width=2.5cm]{./Figuras/Castor.png} & Tres:
			\begin{compactitem}
				\item Rotación alrededor de su eje.
				\item Rotación alrededor de su punto de contacto con el suelo.
				\item Rotación en el eje de la rueda.
			\end{compactitem} & Silla de oficina\\
			\hline
			Omnidireccional (Sueca) \includegraphics[width=2.5cm]{./Figuras/Sueca.png} & Tres:
			\begin{compactitem}
				\item Rotación alrededor de su eje.
				\item Rotación alrededor de su punto de contacto con el suelo.
				\item Rotación alrededor de sus  rodamientos.
			\end{compactitem} & Kuka youBot\\
			\hline
			Esférica \includegraphics[width=2.5cm]{./Figuras/REsfera.png} & Tres:
			\begin{compactitem}
				\item Rotación en cualquier dirección.
				\item Rotación alrededor de su punto de contacto con el suelo.
			\end{compactitem} & Rodamiento de bolas\\
			\hline
		\end{tabular}
	\end{center}
\end{table}
\vspace*{-0.6cm}
\par En el caso de los robots móviles terrestres, los hay con configuraciones basadas en patas (buscando emular el movimiento de los animales, insectos y humanos), orugas y ruedas; de forma similar, dentro de los robots móviles con ruedas se pueden ordenar de acuerdo al tipo de ruedas empleados, y a su vez en la configuración con la que estas funcionan en conjunto para el gobierno de la plataforma. De \cite{barrientossotelovictorricardoRobotsMovilesEvolucion} se ha tomado la Figura \ref{fig:configuraciones}, que ilustra las principales configuraciones de robots móviles, mientras que de \cite{corellnikolausIntroductionAutonomousRobots2014} se tomó y adaptó la Tabla \ref{tab:ruedas} para enumerar los tipos de ruedas máxime usados junto con sus características y usos; la información mostrada en esta Sección se apoya en estas dos bibliografías junto con \cite{fernandoRoboticaControlRobots2011}.
\begin{figure}[htbp!]
	\centering
	\includegraphics[width=0.45\textwidth]{./Figuras/Configuraciones}
	\caption{Principales configuraciones de robots móviles con ruedas.}
	\label{fig:configuraciones}
\end{figure}
\par El desplazamiento de los RMR se da gracias al uso de actuadores, en muchas ocasiones eléctricos. Exceptuando a los motores a pasos (que no resultan convenientes para propulsar un vehículo eléctrico por su baja relación de velocidad), los motores eléctricos de corriente directa e imán permanente son los actuadores más empleados para la locomoción en los RMR, y estos se pueden encontrar a su vez en dos clasificaciones: Con escobillas y sin escobillas.
\par Por motivos prácticos, los motores eléctricos son usados en conjunto con cajas de engranes para reducir su velocidad y aumentar su par; asimismo, al integrar en el sistema un codificador rotatorio\footnote{Dispositivo de naturaleza electromecánica usado para traducir la posición angular de un eje a un código digital.}, es posible conocer la posición y la velocidad del motor para cada instante de tiempo. A estos motores se les conoce como {\it servomotores}, y su control requiere suministrar una serie de pulsos con diferente duración en milisegundos para obtener diferente posición angular; o bien distintos anchos de pulso en milisegundos para distintas velocidades angulares.
\subsection{Odometría}
La odometría es el estudio de la estimación de la posición de los robots móviles durante la navegación. Su estudio implica tomar la información de los sensores presentes en la plataforma para obtener los cambios de posición de la misma a lo largo del tiempo. 
\par El concepto básico es desarrollar un modelo matemático para establecer la manera en que los comandos enviados a las extremidades del robot (ruedas, orugas, patas, etc.) inducen el movimiento de la plataforma. Este término también se usa a veces para referirse a la trayectoria que ha recorrido uno de estos vehículos, dejando la posibilidad de emplear diferentes tipos de sensores para obtener la posición a lo largo del tiempo, como lo son los codificadores rotatorios, sistemas de visión (llamado odometría visual), brújulas, giroscopios, acelerómetros, GPS, etc.
\par La idea fundamental de la odometría se encuentra en la integración de información incremental obtenida a partir del movimiento a lo largo del tiempo, y esto tiene como desventaja una acumulación de errores debido a los propios sensores empleados, en particular la acumulación de errores de orientación (lo que causa grandes inexactitudes en la estimación de la posición), a la larga los errores aumentarán de manera proporcional a la distancia recorrida por el robot. No obstantes estas limitaciones, la odometría es una parte importante del sistema de navegación de un robot, debe usarse junto con medidas del posicionamiento absolutas para proporcionar una estimación de la posición más fiable. Los errores que se mencionan se pueden agrupar en dos categoría:
\paragraph{Errores sistemáticos:}
\begin{itemize}
	\item Los diámetros de las ruedas no son iguales.
	\item Mala alineación de las ruedas.
	\item Baja resolución en la tasa de muestreo de los sensores.
\end{itemize}
\paragraph{Errores no sistemáticos:}
\begin{itemize}
	\item Desplazamiento en suelos desnivelados.
	\item Desplazamiento sobre objetos inesperados que se encuentren en el suelo.
	\item Patinaje de las ruedas (debido a suelos resbaladizos, sobre-aceleración, derrapes, fuerzas externas, falta de un punto de contacto con el suelo).
\end{itemize}
\par Se vuelve un asunto importante distinguir entre los errores sistemáticos y los no sistemáticos para mitigar la incertidumbre en la odometría; en particular los errores sistemáticos son graves, pues se acumulan constantemente; por su parte los errores no sistemáticos generan problemas que pueden aparecer de manera espontánea, y pueden causar errores significativos en la estimación de la posición (un ejemplo se puede dar cuando el robot pasa por encima de un objeto que se encuentra en el suelo).
\subsection{Percepción}
\label{ssec:per}
Es fundamental para cualquier sistema de navegación poseer la habilidad de medir los estados internos del vehículo, además de poder percibir su entorno. Para adquirir el conocimiento del ambiente, los sistemas robóticos emplean diferentes tipos de sensores, y esta actividad se combina con la interpretación de las mismas en función de variables físicas, todo esto bajo la consideración que cada tipo de sensor tiene ventajas y limitaciones inherentes que necesitan ser compensadas. Los sensores involucrados pueden arrojar mediciones de índole externa e interna\footnote{Estas categorías son conocidos como {\it propiocepción} y {\it exterocepción} en el argot biológico.}. El tipo de sensores disponibles para el robot móvil determina de manera directa las estrategias de control disponibles para la navegación autónoma de la plataforma. En \cite{siegwartIntroductionAutonomousMobile2011} y \cite{corellnikolausIntroductionAutonomousRobots2014} se da una descripción más amplia de los principales medios de percepción en la robótica, además se complementa con la tesis \cite{b.VisualHomingCarlike2005}. 
\subsubsection{Sensores internos}
\label{sssec:si}
Esta categoría de sensores es la encargada de la percepción de los estados propios del robot, y abarca la medición del nivel de carga de las baterías, la posición de los actuadores del robot, velocidad, aceleración y torque de los mismos, por sólo mencionar algunas. Entre los sensores propioceptivos de mayor relevancia dentro de los RMR se encuentra el codificador rotatorio: el cual permite conocer la posición angular en los ejes de los motores, y con base en ello se vuelve posible obtener la distancia recorrida y la velocidad de la plataforma, esto al conocer el radio de la rueda y el tiempo entre muestras. El codificador rotatorio, en el caso de un funcionamiento óptico, cuenta con una rueda que tiene ranuras claras y ranuras oscuras, de manera que cuando la rueda se encuentre girando, un emisor óptico se encontrará trabajando de manera continua y por lo tanto el receptor óptico alternará valores lógicos de $1$ y $0$. En la Figura \ref{fig:CRO} se ilustra el funcionamiento de un codificador rotatorio de funcionamiento óptico como el descrito con anterioridad.
\begin{figure}[htbp!]
	\centering
	\includegraphics[width=0.5\textwidth]{./Figuras/Encoder}
	\caption{Disco de codificador rotatorio de funcionamiento óptico.}
	\label{fig:CRO}
\end{figure}
\subsubsection{Sensores externos}
\label{sssec:se} 
Esta jerarquía de sensores es la encargada de la percepción de los aspectos exteriores del robot, y envuelve a la localización de las características del entorno que permiten la ubicación del robot y de los objetos circundantes en el espacio de trabajo. Dentro de esta categoría se encuentran las balizas, las brújulas, los sistemas de posicionamiento global (GPS), los sensores de distancia (ultrasónicos, LIDAR\footnote{Sistema basado en laser que brinda una imagen del entorno a partir de la detección y medición de la distancia alrededor}, etc.), y los sistemas de visión artificial.
\par En esta clase se puede contar con la información provista por los sensores de navegación inercial (denominadas también IMU), que involucran la medición de seis grados de libertad que se dan con la aceleración a lo largo de los ejes $x$, $y$ y $z$, además de los ángulos de \textit{alabeo}, \textit{cabeceo} y \textit{guiñada}. De manera interna, la medición de estas variables se da gracias a los MEMs.
\par Los sensores externos aportan más confianza en las mediciones que los sensores internos porque miden ángulos y distancias de manera directa, lo que resulta más eficiente que hacerlo a partir de la lecturas de velocidad y aceleración. A pesar de ello, estos sensores pueden presentar problemas (mismos que a su vez se pueden mitigar al controlar las condiciones de operación de la plataforma robótica), como lo son:
\begin{itemize}
	\item Requiere una línea de visión para poder operar.
	\item Interferencias por obstrucción o reflexión (en el caso de sensores de visión y de distancia), y por la presencia de estructuras metálicas o campos magnéticos (en el caso de los sensores magnéticos).
\end{itemize}
\subsection{Restricciones holónomas y no-holónomas}
\label{ssec:rnh}
Los robots que usan sólo ruedas con tres grados de libertad son capaces de moverse libremente en el plano. En estos casos, la posición del robot está dada por dos valores que definen la posición y otro que define la orientación. Cada vehículo de ruedas se encuentra sujeto a constantes cinemáticas que lo reducen en su movilidad local, dejando intacta la posibilidad de alcanzar configuraciones arbitrarias con maniobras apropiadas. Si se considera un sistema mecánico cuya configuración $\mathbf{q}$ está descrita por un vector de coordenadas generalizadas, y el espacio de todas las posibles configuraciones del robot coincide con $\mathds{R}^{n}$. El movimiento del sistema que está representado por la evolución de $\mathbf{q}$ en el tiempo, debe estar sujeto a restricciones. Si las mismas pueden ser expresadas en la forma 
\[h_{i}(\mathbf{q})=0\quad i=1, \dots,k<n\]
se les llama holónomas (o integrables). Las restricciones holónomas por lo regular son el resultado de interconexiones mecánicas entre varios elementos del sistema. Por otra parte, las restricciones que involucran coordenadas generalizadas y velocidades en la forma
\[a_{i}(\mathbf{q}, \mathbf{\dot{q}})=0\quad i=1,\dots,k<n\]
son llamadas no-holónomas (también llamadas cinemáticas o no integrables). Las restricciones no-holónomas reducen la movilidad del sistema mecánico, perdiendo accesibilidad en el espacio de configuraciones. En sistemas con restricciones cinemáticas, el número de coordenadas que describe su configuración siempre es mayor que el número de grados de libertad que posee, y en principio ésta es la razón por la que se dificulta estacionar un vehículo en paralelo \cite{sicilianoRoboticsModellingPlanning2009}.
\subsection{Modelado cinemático}
\label{ssec:mc}
Partiendo de la noción de la rueda como un disco rodando en el plano, se puede obtener un modelo generalizado de los RMR a partir del más simple de ellos: El uniciclo.
\par Como se indica en la Figura \ref{fig:Uni}, la posición del robot en el plano se encuentra descrita, con respecto a un marco inercial, por las coordenadas generalizadas $\mathbf{q}=[x, y, \theta]^{T}$, donde el par $(x, y)$ representa las coordenadas de un punto arbitrario sobre el plano sagital\footnote{Plano perpendicular al suelo y en ángulo recto con el plano frontal, que divide al objeto en su mitad derecha e izquierda.} del robot, mientras que $\theta$ describe la orientación del robot con respecto al marco inercial.
\begin{figure}[htbp!]
	\centering
	\includegraphics[width=0.5\textwidth]{./Figuras/Uniciclo}
	\caption{Coordenadas generalizadas para un disco rodando en el plano.}
	\label{fig:Uni}
\end{figure}
\par Para una rueda convencional, las restricciones cinemáticas implican que la velocidad del centro de la rueda es paralela a la plano de rotación ({\it condición de no deslizamiento}) y es proporcional a la velocidad de rotación de la rueda (i. e. {\it condición de balanceo puro}). La condición de balanceo puro para el disco se puede expresar como
\begin{equation}
	\dot{x}\sen\theta-\dot{y}\cos\theta=[\sen\theta, -\cos\theta, 0]\mathbf{\dot{q}}=0
\end{equation}
\par Lo que implica que la velocidad del punto de contacto es cero en la dirección ortogonal al plano sagital de la rueda. La línea que pasa a través del punto de contacto teniendo tal dirección se llama {\it línea de movimiento cero}. Con esto, se puede considerar la matriz
\[\mathbf{G(q)}=[\mathbf{g_{1}(q)}, \mathbf{g_{2}(q)}]=\left[
\begin{array}{cc}
\cos\theta & 0\\
\sen\theta & 0\\
0 & 1\\
\end{array}\right]\]
cuyas columnas $\mathbf{g_{1}(q)}$ y $\mathbf{g_{2}(q)}$ son, para cada $\mathbf{q}$, una base del espacio nulo asociado a las restricciones. Todas las velocidades generalizadas de $\mathbf{q}$ son por lo tanto una combinación lineal de $\mathbf{g_{1}(q)}$ y $\mathbf{g_{2}(q)}$. El modelo cinemático del uniciclo es, entonces
\begin{equation}
	\left[\begin{array}{c}
	\dot{x}\\
	\dot{y}\\
	\dot{\theta}\\
	\end{array}\right]=\left[\begin{array}{c}
	\cos\theta\\
	\sen\theta\\
	0\\
	\end{array}\right]v+
	\left[\begin{array}{c}
	0\\
	0\\
	1\\
	\end{array}\right]\omega
\end{equation}
donde las entradas $v$ y $\omega$ son la velocidad lineal y angular, como corresponde \cite{sicilianoRoboticsModellingPlanning2009}.
\par La configuración de robot móvil que se aproxima más al uniciclo es la configuración diferencial (Figura \ref{fig:configuraciones}c). Entre otras, una de las desventajas de una configuración diferencial es la dependencia de una rueda castor --que se comporta pobremente a altas velocidades y dificulta las estrategias de conducción-- además de requerir de la sincronización de ambos motores para conducirse a la misma velocidad exacta. Por fortuna, las mismas desventajas pueden ser mitigadas por el mecanismo empleado por los autómoviles, que es conducido por un sólo motor y puede gobernar las ruedas laterales a través de un servomecanismo. Esta configuración es conocida como {\it plataforma Ackerman} (Figura \ref{fig:configuraciones}a), donde cada rueda en la configuración Ackerman es su propio punto de pivote y el sistema se encuentra restringido de tal manera que todas las ruedas del vehículo van en una circunferencia con un punto central en común, evitando así derrapes. A bajas velocidades, se puede emplear el modelo de un triciclo o el de una bicicleta para aproximar el comportamiento de la configuración Ackerman. La bicicleta tiene una rueda trasera fija al cuerpo del vehículo, y rota con respecto al eje vertical de la rueda delantera para conducir el vehículo \cite{corellnikolausIntroductionAutonomousRobots2014}.
\par En un robot móvil, si se prolonga con una recta al eje perpendicular a cada rueda, el punto en el que se da la intersección es llamado Centro Instantáneo de Rotación (CIR); este es el punto de referencia sobre el cual el vehículo sigue una trayectoria circular, y entonces su velocidad angular está dada por
\begin{equation}\label{eq:giro}
	\dot{\theta}=\frac{v}{R}
\end{equation}
de manera que por relaciones geométricas se conoce que el radio de giro $R$ se obtiene por $R=L/\tan(\varphi)$, teniendo a $L$ como la distancia entre los ejes frontal y trasero del vehículo, y a $\varphi$ como el ángulo de gobierno para el vehículo, mismo que típicamente se encuentra limitado mecánicamente y su máximo valor determina el valor de $R$. Tomando como ejemplo la Figura \ref{fig:CIR}, las líneas punteadas muestran la dirección a lo largo de la cual las ruedas no se pueden mover, cruzándose en el CIR del vehículo \cite{sicilianoRoboticsModellingPlanning2009, corkeRoboticsVisionControl2017}.
\begin{figure}[htbp!]
	\centering
	\includegraphics[width=0.5\textwidth]{./Figuras/Ackerman}
	\caption{Derecha: Cinemática del mecanismo de gobierno utilizado por los automóviles. Izquierda: Mecanismo de un vehículo Ackerman.}
	\label{fig:CIR}
\end{figure}
\par Gerasimos Rigatos y Krishna Busawon proponen en \cite{rigatosRoboticManipulatorsVehicles2018} un modelo dinámico del vehículo de cuatro ruedas en la configuración Ackerman y otro para un vehículo de dos ruedas para la configuración de bicicleta, además proponen una linealización de estos dos modelos, junto con algunas estrategias de control tanto lineales como no lineales. En dichos modelo se consideran los deslizamientos longitudinal y lateral del vehículo, así como la guiñada, y son representados por las ecuaciones \eqref{eq:long}, \eqref{eq:lat} y \eqref{eq:yaw} respectivamente.
\begin{equation} \label{eq:long}
	-mV(\dot{\beta}+\dot{\varPsi})\sen(\beta)+m\dot{V}\cos(\beta)=f_{x}
\end{equation}
\begin{equation} \label{eq:lat}
	mV(\dot{\beta}+\dot{\varPsi})\cos(\beta)+m\dot{V}\sen(\beta)=f_{y}
\end{equation}
\begin{equation} \label{eq:yaw}
	I\ddot{\varPsi}=T_{z}
\end{equation}
Donde $m$ es la masa del vehículo, $\beta$ representa el ángulo entre el vehículo y el plano sagital, $V$ es el vector velocidad del vehículo, $\varPsi$ es el ángulo de guiñada, $f_{x}$ es la fuerza agregada en el eje $x$, $f_{y}$ es la fuerza agregada en el eje $y$, y $T_{z}$ es el par agregado alrededor del $z$. Estas ecuaciones se pueden escribir de forma matricial como
\begin{equation}
	\left[\begin{array}{ccc}
	-\sen(\beta) & \cos(\beta) & 0\\
	-\cos(\beta) & \sen(\beta) & 0\\
	0 & 0 & 1
	\end{array}\right]
	\left[\begin{array}{c}
	mV(\dot{\beta}+\dot{\varPsi})\\
	m\dot{V}\\
	I\ddot{\varPsi}
	\end{array}\right]=\left[\begin{array}{c}
	f_{x}\\
	f_{y}\\
	T_{z}
	\end{array}\right]
\end{equation}
\par Asimismo, G. Oriolo et al. en \cite{laumondRobotMotionPlanning1998} desarrollan un modelo cinemático para los robots móviles con configuración Ackerman, y el modelo se ha basado a su vez en la bicicleta. Este último trabajo será de más ayuda que el de Gerasimos Rigatos, y se explorará a detalle en la Sección \ref{sec:mc}.
\section{Navegación autónoma}
\label{sec:na}
Algunas de las primeras aproximaciones a la navegación autónoma se remontan al primer par de robots móviles autónomos de la historia. {\it Elmer} y {\it Elsie}\footnote{Siglas en inglés para Electro Light Sensitive Internal External} (Elsie  en la Figura \ref{fig:ELSIE}), un par de tortugas robóticas que fueron construidas a finales de la década de 1940 por William Grey Walter, quien eraun neurocientífico de la universidad de Bristol. Este par de tortugas estaban diseñadas para buscar una fuente de luz sin tener información previa de la posición de la misma \cite{w.greyMachineThatLearns1951}. Durante la década de 1960 apareció el robot {\it Shakey} (Figura \ref{fig:Shakey}) del SRI\footnote{Instituto de Investigación de Stanford}, que era capaz de crear un mapa de su entorno para planear una trayectoria a un punto de destino predefinido \cite{nilsonn.j.ShakeyRobot1984}.
\begin{figure}[htbp!]
	\centering
	\begin{subfigure}[h]{0.4\textwidth}
		\centering
		\includegraphics[width=\textwidth]{./Figuras/ELSIE}
		\caption{Robot Elsie sin su caparazón.}
		\label{fig:ELSIE}
	\end{subfigure}
	\begin{subfigure}[h]{0.4\textwidth}
		\centering
		\includegraphics[width=\textwidth]{./Figuras/Shakey}
		\caption{Robot Shakey.}
		\label{fig:Shakey}
	\end{subfigure}
	\caption{Antecedentes de la robótica móvil}
\end{figure}
\par Una de las funciones más importantes de un robot móvil es poder moverse hacia algún lugar, mismo que puede ser especificado en términos de alguna característica del ambiente (como seguir la luz, seguir paredes, etc.) o en función de alguna coordenada geográfica de referencia. Una estrategia se basa en tomar mediciones simples del mundo y reaccionar a lo medido. Los sistemas de navegación reactiva pueden ser simples y a la vez veloces, ya que el sistema de medición está conectado directamente a la acción. Los sistemas que crean mapas y razonan acerca de estos requieren más recursos; no obstante, son capaces de realizar tareas más complejas \cite{corkeRoboticsVisionControl2017}.
\par Para la robótica móvil, la tarea de navegación representa una ciencia en sí misma por la gran cantidad de técnicas relacionadas con alcanzarla. La localización del robot en el plano, además de la manera en que se ve posibilitado para actuar en éste (esquivando obstáculos, guiándose por diferentes características  del entorno, seguir una ruta, etc.) permiten llevar a la navegación teleoperada o preprogramada a una navegación autónoma.
\par En el contexto urbano actual, se plantean los autos del mañana como sistemas autónomos, i. e. robots móviles autónomos que integren una variedad de sensores para su localización espacial y la consecuente navegación reactiva. Según lo expuesto en \cite{heinrichPlanningUniversalOnRoad2018}, estos sistemas incorporan una gran variedad de beneficios, listados a continuación:
\begin{itemize}
	\item Un vehículo habilitado para conducirse por sí mismo en el escenario y actuar en caso de peligro, puede disminuir el 90 \% de los accidentes de tránsito causados por el ser humano.
	\item Un vehículo habilitado para ir a altas velocidades y cerca de otro sin peligro, reduce el consumo de combustible y emisión de gases a la atmósfera.
	\item Si un vehículo puede moverse de manera autónoma en un terreno, o entrar en un campo minado, o cualquier otra misión peligrosa, el número de individuos en riesgo se decrementa.
\end{itemize}
Además en el mismo texto se sugiere que la operación de los vehículos inteligentes involucra:
\begin{itemize}
	\item El monitoreo del ambiente en torno al vehículo.
	\item Comunicación con la infraestructura de tránsito y otros vehículos.
	\item Percibir la geometría del carril y del camino.
	\item Reconocer los señalamientos de tránsito.
	\item Distinguir el estado de los semáforos.
\end{itemize}
\par Para estas tareas, la visión artificial juega un rol básico. Por un lado, el camino puede ser detectado por sistemas LIDAR (depende de su construcción), sin embargo para detectar los carriles se requiere la visión artificial; la detección de las señales de tránsito se puede realizar tomando en consideración los colores, geometría y reconocimiento de los símbolos presentes dentro de los límites del señalamiento; la detección del estado del semáforo es un problema que involucra reconocer el color y la intensidad del mismo, y por la luz ambiental puede resultar problemático si la misma llega a opacar la luz del semáforo \cite{broggiIntelligentVehicles2008}.
\par La SAE\footnote{Sociedad de Ingenieros Automotrices.} publicó un manual de referencia con los diferentes niveles de automatización \cite{saeTaxonomyDefinitiosTerms2014} para vehículos, el cual contiene las métricas para situar a un vehículo dentro de alguno de los 6 niveles de automatización planteados, mismos que se resumen en la Tabla \ref{tab:niveles}, y van desde la nula automatización (representada por el nivel 0) hasta su antípoda: la nula intervención humana en el proceso de conducción (nivel 5 de esta clasificación).
\begin{center}
	\begin{longtable}{|p{1cm}|p{3.5cm}|p{5.5cm}|p{5.5cm}|}
		\caption{Niveles de automatización propuestos por la SAE}
		\label{tab:niveles}
		\endfirsthead
		\caption*{{\bf Tabla \ref{tab:niveles}:} Continuación}
		\endhead
		\endfoot
		\endlastfoot
		\hline
		{\bf Nivel} & {\bf Designación} & {\bf Rol del conductor} & {\bf Rol de la computadora}\\ 			\hline\hline
		0 & Sin automatización & \begin{compactitem}
			\item Monitoreo del entorno.
			\item Ejecutar las tareas de giro, aceleración y freno.
		\end{compactitem} & \begin{compactitem}
			\item No aplica.
		\end{compactitem}\\ 
			\hline
		1 & Conducción asistida & \begin{compactitem}
			\item Monitoreo del entorno.
			\item Realizar el control lateral o longitudinal según corresponda.
			\item Desactivar la conducción asistida si se requiere.
		\end{compactitem} & \begin{compactitem}
			\item Realizar el control lateral o longitudinal según corresponda.
		\end{compactitem}\\
			\hline
		2 & Automatización parcial & \begin{compactitem}
			\item Monitoreo del entorno.
			\item Determinar el momento adecuado para la activación de la automatización parcial.
			\item Supervisar la conducción del sistema de automatización parcial.
			\item Desactivar la automatización parcial si se requiere.
		\end{compactitem} & \begin{compactitem}
			\item Ejecutar las tareas de giro, aceleración y freno.
		\end{compactitem}\\
			\hline
		3 & Automatización condicional & \begin{compactitem}
			\item Determinar el momento adecuado para la activación de la automatización condicional.
			\item Supervisar la conducción del sistema de automatización condicional.
			\item Desactivar la automatización condicional si se requiere.
		\end{compactitem} & \begin{compactitem}
			\item Monitoreo del entorno.
			\item Ejecutar las tareas de giro, aceleración y freno.
		\end{compactitem}\\
			\hline
		4 & Automatización alta & \begin{compactitem}
			\item Determinar el momento adecuado para la activación de la automatización condicional.
			\item Supervisar la conducción del sistema de automatización condicional.
			\item Desactivar la automatización condicional si se requiere.
		\end{compactitem} & \begin{compactitem}
			\item Monitoreo del entorno.
			\item Ejecutar las tareas de giro, aceleración y freno.
			\item Permitir la activación sólo bajo los casos para los cuales ha sido diseñado.
			\item Desactivarse hasta que el conductor humano tome el control.
		\end{compactitem}\\
			\hline
		5 & Automatización completa & \begin{compactitem}
			\item No aplica.
		\end{compactitem} & \begin{compactitem}
			\item Monitoreo del entorno.
			\item Ejecutar las tareas de giro, aceleración y freno.
		\end{compactitem}\\
			\hline
		\end{longtable}
	\end{center}
\section{Visión}
\label{sec:vi}
Con muy pocas y contadas excepciones, todos los animales cuentan con ojos, y estos están asociados a un poderoso (e igualmente complejo) sistema de procesamiento paralelo para el volumen de información percibida, mismo sistema que ha tomado millones de años de constante evolución para alcanzar su estado de desarrollo actual. La visión es una poderosa habilidad para la navegación en los animales y en seres humanos, y es por ello que se posiciona como una buena elección para la adquisición de información cualitativa que permite la navegación de los robots móviles. Los sistemas de visión forman parte de los sensores exteroceptivos, pero al abarcar por sí mismo un amplio campo de investigación, en esta Sección se describirá por separado de lo expuesto en la \autoref{ssec:per}. La redacción de esta Sección se apoya sobre todo haciendo referencia a Peter Corke, Richard Szeliski y a Heredia Favieri en \cite{corkeRoboticsVisionControl2017},  \cite{szeliskiComputerVisionAlgorithms2011} y \cite{herediafavieriReconocimientoObjetosImagenes2015}, como corresponde.
\begin{figure}[htbp!]
	\centering
	\includegraphics[width=0.7\textwidth]{./Figuras/Espectro}
	\caption{Colores del espectro visible en función de su longitud de onda.}
	\label{fig:espectro}
\end{figure}
\par Fue hacia el año 1670 cuando Sir Isaac Newton se dio cuenta de que la luz blanca es el resultado de la combinación de luces de diferentes colores. Ahora se sabe que cada uno de esos colores está asociado con una longitud de onda y a una frecuencia en el espectro electromagnético, pudiendo observar nosotros como humanos un pequeño rango de éste espectro, abarcando longitudes de onda que abarcan de manera aproximada entre 400 y 700 nm aproximadamente, como se puede observar en la Figura \ref{fig:espectro}.
\begin{figure}[htbp!]
	\centering
	\begin{subfigure}[h]{0.25\textwidth}
		\centering
		\includegraphics[width=\textwidth]{./Figuras/RGB}
		\caption{Modelo RGB.}
		\label{fig:RGB}
	\end{subfigure}
	\begin{subfigure}[h]{0.25\textwidth}
		\includegraphics[width=\textwidth]{./Figuras/HSV}
		\caption{Modelo HSV.}
		\label{fig:HSV}
	\end{subfigure}
	\begin{subfigure}[h]{0.25\textwidth}
		\includegraphics[width=\textwidth]{./Figuras/CIE}
		\caption{Modelo CIE.}
		\label{fig:CIE}
	\end{subfigure}
	\caption{Algunos de espacios de color más empleados.}
\end{figure}
\par El modelo más empleado en la actualidad para la representación de los colores es el que utiliza como base al Rojo, Verde y Azul (RGB por sus siglas en inglés), y es el más natural al aproximarse a la manera en la que el ojo humano capta la luz como colores positivos y negativos. En la Figura \ref{fig:RGB} se muestra la manera en que al combinar estos colores entre sí en diferentes intensidades se obtiene como resultado otro color diferente, siendo el blanco la combinación de los tres en la máxima intensidad, y el negro la ausencia de luz. Sin embargo, el espacio de color RGB no es el único que puede ser utilizado para representar los colores en la luz, y algunos de estos ejemplos son el HSV y el CIE (Figuras \ref{fig:HSV} y \ref{fig:CIE}.), aunque se debe decir que existen otros tantos y además variaciones de los mismos.
\begin{figure}[htbp!]
	\centering
	\includegraphics[width=0.5\textwidth]{./Figuras/Bayer}
	\caption{Patrón de Bayer para una cámara.}
	\label{fig:Bayer}
\end{figure}
\par Las imágenes, computacionalmente hablando, se componen de la intensidad lumínica que se percibe en cada uno de los elementos de un arreglo bidimensional de receptores sensibles a una longitud de onda específica. La información provista por la imagen equivale a la captura de una escena en un instante, y representa los valores de la intensidad que llega a cada uno de los receptores lumínicos de la cámara. En una cámara digital que trabaja en el espacio de color RGB, por ejemplo, se cuenta con un patrón de receptores lumínicos llamado patrón de Bayer (Figura \ref{fig:Bayer}), y estos receptores pueden existir como semiconductores con tecnología CCD o CMOS.
\begin{figure}[htbp!]
	\centering
	\includegraphics[width=0.6\textwidth]{./Figuras/Perspectiva}
	\caption{Proyección de perspectiva tridimenional a bidimensional.}
	\label{fig:Per}
\end{figure}
\par Gracias a la información obtenida por las imágenes, es posible conocer características importantes de los objeto que se encuentran en el escenario capturado por la cámara, tales como el color, la forma, la posición o la orientación, por mencionar algunas. El proceso de formación de imágenes en un ojo o una cámara involucra la proyección de un mundo tridimensional en una superficie bidimensional. La información de profundidad está perdida cuando se trabaja sólo con una cámara, por lo que no se puede distinguir al punto más cercano del más lejano en una imagen. La transformación de imágenes desde un espacio tridimenensional a uno bidimensional se conoce como proyección de perspectiva, y en la Figura \ref{fig:Per} se muestra un diagrama que expone la naturaleza de la transformación, en la misma el cuadro sombreado con una línea verde punteada rodeándolo representa la imagen bidimensional, en donde un punto $P=(X, Y, Z)$ adquiere coordenadas $p=(u, v)$, pudiendo percibir de esta manera la pérdida de la información de profundidad en la proyección de perspectiva.
\par Además de la pérdida de información provocada por la proyección de perspectiva, surge una incertifumbre cuando la imagen que se observa sufre deformaciones debido a diferentes factores, como la ubicación de la cámara con respecto al objeto a observar. Ante todo existen dos acciones por las que se pueden producir alteraciones en la imagen:
\begin{enumerate}
	\item La interpolación de los niveles intensidad de los pixeles.
	\item Una transformación espacial que define la reubicación de los pixeles en la imagen.
\end{enumerate}
\par La primera de estas acciones, la interpolación de los niveles de intensidad en los pixeles, es un problema que puede originarse por problemas de la iluminación de la escena o del objeto, mientras que por otra parte puede surgir por la sensibilidad de los receptores. En realidad estos problemas en la interpolación pueden ser resueltos al controlar la manera en que se toma la captura de la escena, buscando controlar la iluminación de acuerdo a las características que se busca capturar.
\par En la segunda de las acciones, que implica la perspectiva de la cámara con respecto a la escena, se involucran a las transformaciones espaciales. Estas también llamadas modificaciones, preservan los ángulos mientras que las distancias entre los puntos cambian en una proporción fija y los objetos mantienen una forma similar (triángulos similares, círculos similares, etc.). Estas transformaciones contienen el escalamiento, la traslación y la rotación. De manera general se usa la operación matricial de la Ecuación \eqref{eq:simi} para calcular las transformaciones de similaridad:
\begin{equation}\label{eq:simi}
	\left[\begin{array}{c}
	x'\\
	y'\\
	1
	\end{array}\right]=
	\left[\begin{array}{ccc}
	a\cos\theta & -a\sen\theta & x_{t}\\
	a\sen\theta & a\cos\theta & y_{t}\\
	0 & 0 & 1
	\end{array}\right]
	\left[\begin{array}{c}
	x\\
	y\\
	1
	\end{array}\right]
\end{equation}
\par Si a las transformaciones de similaridad las acompañan las transformaciones de deslizamiento y escalado no uniforme, se obtienen las transformaciones afines, de las cuales se destaca que los objetos que son observados conservan las líneas paralelas. La operación que definen las transformaciones afines se hayan en la Ecuación \eqref{eq:afin}, conteniendo los componentes de seis grados de libertad correspondientes a los coeficientes $a_{00}$, $a_{01}$, $a_{10}$, $a_{11}$.
\begin{equation}\label{eq:afin}
	\left[\begin{array}{c}
	x'\\
	y'\\
	1
	\end{array}\right]=
	\left[\begin{array}{ccc}
	a_{00} & a_{01} & x_{t}\\
	a_{10} & a_{11} & y_{t}\\
	0 & 0 & 1
	\end{array}\right]
	\left[\begin{array}{c}
	x\\
	y\\
	1
	\end{array}\right]
\end{equation}
\par Junto con la identificación de las transformaciones espaciales que sufre la imagen, existen algunas otras operaciones que son de importancia y a la vez muy comunes dentro del procesamiento de imágenes. Se toman como elementos al filtrado de imágenes y la detección de sus características.
\par Al ser las imágenes señales bidimensionales, la teoría de filtros sobre señales unidimensionales puede ser extendida para describir el funcionamiento de los filtros en dos dimensiones. Existen filtros tipo pasa alto, pasa bajo, pasa banda y rechaza banda por mencionar sólo algunos. El proceso de filtrado puede llevarse a cabo en el dominio de la frecuencia o en el dominio del espacio. Por tanto, se pueden definir los principales objetivos que se consiguen con el filtrado de imágenes:
\begin{itemize}
	\item {\bf Suavizar la imagen:} Reducir la cantidad de variaciones de intensidad entre pixeles vecinos.
	\item {\bf Eliminar ruido:} Eliminar aquellos pixeles cuyo nivel de intensidad sea muy diferente al de sus vecinos y cuyo origen pueda estar tanto en el proceso de adquisición de la imagen como en el de transmisión.
	\item {\bf Realzar bordes:} Destacar los bordes que se localizan en una imagen.
	\item {\bf Detectar bordes:} Detectar los pixeles donde se produce un cambio brusco de la intensidad.
\end{itemize}
\par Si se requiere realizar el filtrado de la imagen en el dominio de la frecuencia, entonces se emplea la transformada de Fourier sobre la imagen. De manera general, el procedimiento empleado para aplicar los filtros en el dominio de la frecuencia a una imagen es el siguiente:
\begin{enumerate}
	\item Se aplica la transformada de Fourier (FFT\footnote{Transformada Rápida de Fourier.}).
	\item El resultado se multiplica por la función del filtro que ha sido escogido (Butterworth, Bessel, Chebyshev, etc.).
	\item Se aplica la transformada inversa de Fourier (IFFT\footnote{Transformada Inversa Rápida de Fourier.}) para devolver la imagen al dominio espacial.
\end{enumerate}
\par Mientras que para realizar el filtrado en el dominio del espacio, las operaciones se llevan a cabo de manera directa sobre los pixeles de la imagen, esto con el uso de núcleos de dimensión cuadrada cuyo objetivo es resaltar los cambios en la imagen. Durante el proceso se relaciona, para cada uno de los pixeles que conforman la imagen, un conjunto de pixeles próximos al pixel objetivo con la finalidad de obtener información útil relacionado con el tipo de filtro que se ha aplicado, lo cual permite actuar sobre el pixel concreto en que se está llevando a cabo el proceso de filtrado. A través del filtrado en el dominio del espacio se puede trabajar sobre la detección de características de imágenes como los bordes, siendo algunos de los más conocidos núcleos u operadores el Sobel, Prewitt, Roberts o Canny.
\par Adentrándose en las múltiples aplicaciones de Visión Artificial, existen múltiples operaciones que resultan útiles para el análisis de imágenes, y la detección de características. De la documentación de la librería OpenCV\cite{opencvOpenCVReferenceManual} se toma la Tabla \ref{tab:car}, donde se exponen algunos de los principales algoritmos para detección de características en imágenes.
\begin{center}
	\begin{longtable}{|l|l|}
		\caption{Principales algoritmos para detección de características empleados por OpenCV.}
		\label{tab:car}
		\endfirsthead
		\caption*{{\bf Tabla \ref{tab:car}:} Continuación.}
		\endhead
		\endfoot
		\endlastfoot
		\hline
		{\bf Tipo de operador} & {\bf Nombre del algoritmo}\\
		\hline
		\hline
		\multirow{5}{*}{}
		& Sobel\\ \cline{2-2}
		& Prewitt\\ \cline{2-2} 
		Detección de bordes & Deriche\\ \cline{2-2}
		& Roberts\\ \cline{2-2}
		& Canny\\ \hline
		\multirow{5}{*}{}
		& Harris\\ \cline{2-2}
		& Shi y Tomasi\\ \cline{2-2}
		Detección de esquinas & Curvas de nivel\\ \cline{2-2}
		& SUSAN\\ \cline{2-2}
		& FAST\\ \hline
		\multirow{4}{*}{}
		& Laplaciano del Gaussiano (LoG)\\ \cline{2-2}
		& Diferencia del Gaussiano (DoG)\\ \cline{2-2}
		Reconocimiento de regiones & Determinante del Hessiano (DoH)\\ \cline{2-2}
		& PCBR\\ \hline
		\multirow{2}{*}{}
		Transformada de Hough & Generalizada\\ \cline{2-2}
		& Probabilística\\ \hline
		\multirow{2}{*}{}
		Tensor estructural & Generalizado\\ \cline{2-2}
		& Probabilístico\\ \hline
		\multirow{3}{*}{}
		& Adaptación de figura afín\\ \cline{2-2}
		Detección de características afines invariantes & Harris afín\\ \cline{2-2}
		& Hessiano afín\\ \hline
		\multirow{5}{*}{}
		& SIFT\\ \cline{2-2}
		& SURF\\ \cline{2-2}
		Descripción de características & GLOH\\ \cline{2-2}
		& HOG\\ \hline
	\end{longtable}
\end{center}
\subsection{Visión tridimensional}
El mundo se compone de elementos tridimensionales y dinámicos, en él los hechos acontecen a alrededor en las dimensiones horizontal, vertical y en profundidad. Una imagen bidimensional aislada no da información suficiente en cuanto a la profundidad de los objetos en una escena, por lo que a grandes rasgos sería necesario definir la generación de un mundo tridimensional a través de tres componentes: el mundo de objetos tridimensionales, las fuentes de luz que componen y la cámara que observa la escena. De forma tradicional, durante gran parte del siglo XX y algunos años del siglo XXI se empleó la combinación de multiples cámaras para el cálculo de puntos tridimensionales a partir de las perspectivas individuales de cada una de las cámaras involucradas. Buscando de este modo emular el funcionamiento de la visión humana, se construyeron cámaras binoculares con disposición similar al acomodo de los ojos,  llamadas cámaras estereoscópicas. Un enfoque más moderno se da con la incorporación de los sensores RGB-D, que involucran la combinación de una cámara en el espacio de color RGB junto con un sensor de profundidad en la escena. Los siguientes dos apartados abarcan estos dos enfoques de visión tridimensional.
\subsubsection{Visión estéreo}
El concepto básico que se encuentra latente en la visión estéreo es la triangulación: Desde la ubicación de un punto de la escena, se forma un triángulo entre la correspondecia del mismo punto en las imágenes de las dos cámaras que toman la escena; conociendo la distancia existente entre las dos cámaras, además de los ángulos formados entre las cámaras y el punto de la escena, la distancia hasta el objeto puede ser determinada. La extracción de la estructura tridimensional de una escena a partir de imágenes estereoscópicas es un problema que ha sido estudiado ampliamente.
\par El funcionamiento esencial de las cámaras estereoscópicas tiene lugar en la geometría epipolar, la cual sólo tiene influencia de los parámetros intrínsecos de la cámara utilizada y su posición relativa, y es independiente de la estructura de la escena. Tal como se observa en la Figura \ref{fig:epi}, la escena de una cámara no es más que una traslación y una rotación de la imagen original.
\begin{figure}[htbp!]
	\centering
	\includegraphics[width=0.6\textwidth]{./Figuras/Epipolar}
	\caption{Modelo de geometría epipolar para cámara estereoscópica.}
	\label{fig:epi}
\end{figure}
\par En este punto inicia por dilucidarse la triangulación que es necesaria aplicar para obtener las coordenadas espaciales de un punto de la escena observado por ambas cámaras, esto a partir de las coordenadas planares de ambas cámaras. La desventaja principal de la visión estereoscópica se encuentra en el poder de cómputo necesario para poder manejar ambas imágenes como una sola y además realizar las operaciones de procesamiento necesarias. A pesar de esto, su uso es ampliamente extendido en general con procesos donde se requiere reducir costos.
\subsubsection{Sensores RGB-D}
Son muchos los sensores que se desarrollan en primer lugar para distintas aplicaciones diferentes antes de llegar a usarse en la robótica, tal es el caso de los sensores RGB-D, mismos que fueron desarrollados para la industria militar antes de pasar a la industria de los videojuegos (como el caso del Kinect de Microsoft), para de ahí expandirse a una amplia variedad de campos, incluyendo sus aplicaciones a la robótica. El Kinect (Figura \ref{fig:rgbd}) es un dispositivo que emplea una cámara como sensor RGB, además de un proyector de haces de luz infrarroja y un receptor (estos dos integran el sensor de profundidad). Las cámaras RGB-D proveen información de importancia acerca del entorno, pues a través de una nube de puntos y una imagen plana se vuelve posible tener un panorama tridimensional de alrededor; similar a lo que pudiera ser obtenido a partir de cámaras estéreo, con un costo de cómputo menor.
\begin{figure}[htbp!]
	\centering
	\includegraphics[width=0.6\textwidth]{./Figuras/RGB-D}
	\caption{Esquema de la configuración interna de un sensor RGB-D Kinect.}
	\label{fig:rgbd}
\end{figure}
\par Para el reconocimiento de objetos por medio de este tipo de sensores se emplean tres etapas: Extracción de puntos de interés, generación de descriptores y búsqueda de correspondencias. En general, la generación de descriptores se refiere a la extracción de representaciones significativas de los objetos a partir de observaciones como mallas poligonales o nubes de puntos. La imagen ya se ha descrito a lo largo de esta sección como una señal bidimensional, sin embargo quedan por definir las mallas poligonales y las nubes de puntos, mismas relaciones que son útiles para expresar el funcionamiento de los sensores RGB-D.
\paragraph{Malla poligonal:} Es una colección de vértices, ejes y caras que definen la forma de un objeto en tres dimensiones. Un vértice corresponde a un punto en el espacio, un eje es una conexión entre dos vértices, y una cara es un conjunto cerrado de ejes, siendo en general un triángulo o cuadrilátero. La densidad de una malla poligonal determina la resolución del objeto.
\paragraph{Nube de puntos:} Es una colección de puntos tridimensionales. Puede ser generada de forma artificial o provenir de capturas de elementos del mundo. Las coordenadas $\{x, y, z\}$ de cualquier punto de la nube están dadas con respecto a un sistema de coordenadas fijo. Si la nube representa datos del mundo, entonces el origen del sistema de coordenadas suele ser el sistema de captura utilizado. El valor de cada punto representa la distancia desde el origen hasta la superficie donde el punto fue capturado, pudiendo incluirse en cada punto datos como el color o la intensidad lumínica para esa pequeña sección de superficie en el objeto.
\par Para mostrar las diferencias en cuanto a su representación, en la Figura \ref{fig:elef} se muestra una malla poligonal y una nube de puntos, ambas del mismo objeto.
\begin{figure}[htbp!]
	\centering
	\includegraphics[width=0.5\textwidth]{./Figuras/Elefante}
	\caption{Derecha: Malla poligonal. Izquierda: Nube de puntos.}
	\label{fig:elef}
\end{figure}
\section{Control}
\label{sec:con}
El control consiste en manipular el valor de una variable para que el comportamiento se mantenga en un valor de referencia. En la actualidad, la teoría del control ha arrojado múltiples enfoques de control como la lógica difusa, redes neuronales, control clásico, control moderno o control robusto; sin embargo para llevar a cabo el control de un vehículo autónomo se considerarán tres enfoques:
\begin{itemize}
	\item Enfoque con Inteligencia Artificial.
	\item Enfoque sin Inteligencia Artificial (ingeniería manual).
\end{itemize}
\subsection{Enfoque con Inteligencia artificial}
La inteligencia artificial representa a la inteligencia llevada a cabo por máquinas. El término tal cual se aplica a las máquinas que son capaces de imitar las funciones cognitivas relacionadas con: Percibir, razonar, aprender y resolver problemas. La inteligencia artificial tiene aplicaciones en áreas como el control de sistemas, planificación automática, reconocimiento de patrones, etc. La máquina biológica que es capaz de realizar estas acciones es el cerebro a través de las neuronas, a partir de lo cual han surgido aproximaciones basadas en las mismas, y son llamadas neuronas artificiales. El agrupamiento de las neuronas --naturales o artificiales-- compone una red neuronal. La Figura \ref{fig:RNN} muestra la estructura general de una neurona natural en comparación de una neurona artificial con la Figura \ref{fig:RNA}.
\begin{figure}[htbp!]
	\centering
	\begin{subfigure}[h]{0.4\textwidth}
		\centering
		\includegraphics[width=\textwidth]{./Figuras/RNN}
		\caption{Representación esquemática de una neurona natural.}
		\label{fig:RNN}
	\end{subfigure}
	\begin{subfigure}[h]{0.5\textwidth}
		\centering
		\includegraphics[width=\textwidth]{./Figuras/RNA}
		\caption{Representación esquemática de una neurona artificial.}
		\label{fig:RNA}
	\end{subfigure}
\end{figure}
\par Mediante el enfoque basado en IA no es necesario calcular la velocidad del vehículo ni el ángulo de giro en las ruedas usando ecuaciones matemáticas. Se propone una solución al problema de múltiples tareas para el gobierno del vehículo (lo que incluiría la evasión de obstáculos y el control de la orientación), y este enfoque permitiría llevar al vehículo desde su configuración inicial hasta su configuración final; es apoyado por el uso de redes neuronales, devolviendo como resultado el ángulo para la dirección de las ruedas y la velocidad del vehículo. Se emplea un agente inteligente\footnote{En Inteligencia Artficial, es una entidad capaz de percibir su entorno, procesar tales percepciones y responder o actuar en su entorno de manera racional.} que elige la mejor acción de acuerdo a los datos provistos por los sensores para la coordinación de los movimientos del robot. 
\subsection{Enfoque sin Inteligencia Artificial}
El enfoque sin Inteligencia Artificial usa el modelo matemático del vehículo para calcular la dirección de las ruedas necesaria para mantener el vehículo en el rumbo deseado, que es detectado a través de la visión artificial para las características del entorno. Uno de los métodos más populares en la teoría de control es el controlador PID (Proporcional-Integral-Derivativo). El controlador trabaja en un bucle que de manera continua calcula el valor de un error $e(t)$ como una diferencia entre la siguiente la orientación deseada y la actual (para este caso). A partir de esto la compensación es calculada y aplicada. El valor de correción $u(t)$ consiste en tres partes (Proporcional, Integral y Derivativa) y puede ser calculado a partir del error $e(t)$, como se aprecia en la Figura \ref{fig:PID}.
\begin{figure}[htbp!]
	\centering
	\includegraphics[width=0.5\textwidth]{./Figuras/PID}
	\caption{Diagrama de bloques para un controlador PID.}
	\label{fig:PID}
\end{figure}
\par En la expresión matemática para un controlador PID cada coeficiente ($K_{p}$, $K_{i}$, $K_{d}$) necesita ser sintonizado. Ahora, la expresión matemática para una controlador PID es la siguiente:
\[u(t)=K_{p}e(t)+K_{i}\int_{0}^{t}e(\tau)d\tau+K_{d}\frac{de(t)}{dt}\]
\par Para el problema que se aborda en este trabajo también se han encontrado en la bibliografía enfoques usando {\it Control Predictivo Basado en Modelo (MPC)} o {\it Control Óptimo}, principalmente. 
\subsection{Control con realimentación visual}
Este trabajo se centra en la navegación reactiva a través de ambientes desconocidos, por lo que una vez obtenidas las directivas por parte del sistema de visión, es necesario pasarlas al sistema de control; es una sinergia que en la bibliografía se aborda como {\it visual servoing}, sin embargo el mismo se encuentra documentado en su mayoría para su aplicación en manipuladores robóticos y en comparación son pocos los textos en los que se aborda el problema para la robótica móvil. Partiendo de la metodología empleada por estos últimos es que se desarrolla esta sección.
\par En términos generales, el control con realimentación visual implica tomar la información obtenida de los sistemas visuales presentes en la plataforma, y usarla para controlar el movimiento del robot. De manera particular, existen dos técnicas principales que se toman en cuenta al momento de realizar el control visual: IBVS\footnote{Control visual basado en imagen.} y PBVS\footnote{Control visual basado en posición.}. 
\subsubsection{IBVS}
Se trata de un enfoque de control aplicado a la robótica, está basado en tareas definidas visualmente en coordenadas cartesianas con las características del entorno. A partir de esto se pueden definir los errores de aproximación y orientación para el IBVS como
\begin{subequations}
	\begin{equation}\label{eq:eibvs}
		e_{IBVS}=\sqrt{\dot{x}^{2}+\dot{y}^2}
	\end{equation}
	\begin{equation}\label{eq:dibvs}
		\delta_{IBVS}=\arctan\left(\frac{\dot{y}}{\dot{x}}\right)
	\end{equation}
\end{subequations} 
teniendo a $\dot{x}$  y a $\dot{y}$ como las componente en $x$ e $y$ de la velocidad instantánea del RMR.
\begin{figure}[htbp!]
	\centering
	\includegraphics[width=0.5\textwidth]{./Figuras/IBVS/IBVS}
	\caption{Control visual basado en IBVS.}
	\label{fig:IBVS}
\end{figure}
\par A partir de estos dos parámetros es suficiente para estimar la ubicación del robot de manera planar. En el control basado en IBVS las características de la imagen son usadas directamente como realimentación en el sistema, con lo cual puede reducir el retraso en el cálculo, además de eliminar la necesidad de interpretación de la imagen y eliminar errores causados por la calibración de la imagen. Sin embargo representa por sí mismo un reto para el diseño del controlador, esto debido a que se cuenta con una proceso no lineal. Un diagrama de bloques que representa al IBVS se puede encontrar en la Figura \ref{fig:IBVS}.
\subsubsection{PBVS}
El diagrama de bloques para un sistema de control basado en PBVS se observa en la Figura \ref{fig:PBVS}. En el control basado en PBVS las características son extraídas de la imagen y son usadas en conjunto con el modelo geométrico del objetivo y en la posición estimada de la cámara. La realimentación es calculada al reducir el error en el espacio de configuraciones estimado.
En el control basado en PBVS los errores se calculan de manera planar como
\begin{subequations}
	\begin{equation}\label{eq:epbvs}
		e_{PBVS}=\sqrt{\dot{x}^{2}+\dot{y}^2}
	\end{equation}
	\begin{equation}\label{eq:dpbvs}
		\delta_{PBVS}=\theta_{d}-\theta
	\end{equation}
\end{subequations}
teniendo a $\dot{x}$ e $\dot{y}$ como las componentes $x$ e $y$ de la velocidad instantánea, respectivamente, $\theta_{d}$ como la orientación deseada y $\theta$ como la orientación actual.
\par Los componentes $e_{IBVS}$, $\delta_{IBVS}$, $e_{PBVS}$ y $\delta_{PBVS}$ deben pertenecer a la interacción del plano de la imagen con el espacio de trabajo. A esa interacción se le conoce como matriz Jacobiana, y su resultado son los componentes $\dot{x}$ e $\dot{y}$.
\begin{figure}[htbp!]
	\centering
	\includegraphics[width=0.5\textwidth]{./Figuras/PBVS/PBVS}
	\caption{Control visual basado en PBVS.}
	\label{fig:PBVS}
\end{figure}
	\chapter{Estado del arte}
	\label{cap:eda}
	En este Capítulo se engloban algunos de los antecedentes importantes del desarrollo de la conducción autónoma en los robots móviles. Los trabajos tomados en cuenta no llevan un orden cronológico, si no que su contenido es semántico, separando los antecedentes similares y mencionándolos de manera corrida sin realizar mayor división.
\par Este trabajo parte de los llamados vehículos autónomos a escala 1:10, en particular el {\it AutoNOMOS Model Car} (Figura \ref{fig:auto}), donados por el Dr. Raúl Rojas y la {\it Freie Universität Berlin}. Similares a este modelo, también se encuentran en desarrollo vehículos de esta clase por parte de como el {\it MIT} \cite{okellyF110OpenSource2019}, el {\it Worcester Polytechnic Institute} \cite{chenSelfdriving10Race2018}, por sólo mencionar algunas de las universidades que han desarrollado vehículos similares. Los vehículos incorporan una variedad de componentes que los distinguen entre ellos, como variaciones en los sensores de adquisición de imágenes (cámaras RGB-D, camaras estereoscópicas, cámaras omnidireccionales y cámaras monoculares), las unidades de procesamiento (basadas en CPU como la Raspberry Pi u Odroid, y las basadas en GPU).
\begin{figure}[htbp!]
	\centering
	\includegraphics[width=0.5\textwidth]{./Figuras/Mini}
	\caption{AutoNOMOS Model.}
	\label{fig:auto}
\end{figure}
\par Los estudios desarrollados con estos vehículos abarcan el control de los mismos a través de distintas técnicas; tomando como ejemplo el trabajo de Juan Manuel Ibarraza et al.  \cite{juanmanuelibarrazannathaGeneracionTrayectoriasPara}, donde se investigó la generación de trayectorias para esta clase de vehículos mediante el control de modos deslizantes. En \cite{StreetMarkDetection} se realiza la detección de marcas del camino con una cámara web en combinación con una microcomputadora Raspberry Pi 2 a bordo de un vehículo a control remoto a escala 1:10; el vehículo está acondicionado con la computadora y la webcam a boro. La cámara trabaja en el espacio de color HSV, utilizando la librería OpenCV sobre la microcomputadora para realizar el procesamiento de la imagen, el cual emplea un filtrado gaussiano para el suavizamiento de la imagen.
\par Es importante mencionar al gran conjunto de investigaciones en torno a la conducción autónoma en vehículos de tamaño real, de donde han venido gran parte de las aportaciones. El DARPA Grand Challenge es de especial importancia al haber aportado vehículos como el Stanley o el Spirit o Berlin, cuyas aportaciones en el campo son la principal influencia para desarrollos más recientes como el Tesla Autopilot o el Google Waymo. A propósito del Spirit of Berlin, se rescata en \cite{javierrojoSpiritBerlinAutonomous2017} el informe técnico en el que se detallan los procedimientos seguidos por el equipo para la realización de este vehículo autónomo, al que llamaron {\it Spirit of Berlin} (Figura \ref{fig:Spirit}).
\begin{figure}[htbp!]
	\centering
	\includegraphics[width=0.5\textwidth]{./Figuras/SoB}
	\caption{Vehículo Spirit of Berlin, participante en el Darpa Grand Challenge.}
	\label{fig:Spirit}
\end{figure}
\par Las exploraciones alrededor vehículos autónomos, que se desencadenó con lo ya mencionado, ha llevado al desarrollo de nuevas técnicas de control enfocadas a esto; Katrakazas et al. mencionan en \cite{katrakazasRealtimeMotionPlanning2015} una perspectiva técnica histórica a través de varias investigaciones. Dentro del objetivo de buscar la autonomía, un grupo de investigadores chinos se encargó de desarrollar en \cite{wangLateralControlAutonomous2015} un controlador lateral para vehículos autónomos basado en lógica difusa. El sistema se basa en un controlador difuso tipo Mamdani y consiste en realizar el análisis del comportamiento de una trayectoria predefinida como referencia. La velocidad es considerada como un antecedente para las reglas de control difusas. Además de las simulaciones por computadora, los resultados se probaron con robots diferenciales del modelo IN$^{2}$BOT.
\par En \cite{marinoNestedPIDSteering2011}, Riccardo Marino et al. exploran las técnicas de preservación del carril para vehículos autónomos con base en un controlador PID no lineal con realimentación visual. En el modelo propuesto, la dinámica longitudinal se encuentra desacoplada de la dinámica lateral y puede ser despreciado para propósitos de control, lo cual permite linealizar el modelo. Esta investigación se simuló computacionalmente en el entorno CarSim, comparando los resultados con la respuesta de un conductor real. 
\par Sobre una investigación similar a la anterior, en \cite{leeSynthesisRobustLane2018} se aborda la preservación del carril, donde la planta es la dinámica lateral del vehículo, y el control tiene la característica de ser predictivo a pesar de las alteraciones exteriores que pudieran afectar la entrada. La investigación compara los resultados de simulación computacional entre un controlador PID, un regulador lineal cuadrático gaussiano, y un control tipo H$_{\infty}$. También en \cite{mackunisUnifiedTrackingRegulation2014}, se emplea una cámara monocular como sensor en un robot no-holónomo. El control se basa en el modelo geométrico de la perspectiva de la cámara a medida que el robot avanza, todo en el esquema IBVS.
\par Junto con la propia navegación autónoma surge la necesidad de contar con escenarios que faciliten este proceso a través de la propia comunicación de los vehículos con el entorno y viceversa; al respecto existen elementos importantes como AUTOPÍA, desarrollado en la Universidad Politécnica de Madrid en colaboración con la Universidad Complutense de Madrir, y esto es abordado por Pérez Rastelli en su tesis \cite{joshuemanuelperezrastelliAgentesControlVehiculos2012}, de donde se rescata la imagend e la Figura \ref{fig:Autopia}.
\begin{figure}[htbp!]
	\centering
	\includegraphics[width=0.5\textwidth]{./Figuras/AUTOPIA}
	\caption{Pista principal del proyecto AUTOPÍA.}
	\label{fig:Autopia}
\end{figure}
	\chapter{Planteamiento del problema}
	\label{cap:pdp}
	Como Trabajo Terminal, se construirá un robot móvil autónomo, con la capacidad de reconocer el carril por el que se desplaza mediante visión artificial; como lo haría un vehículo real que es operado por un humano. Para validar este robot móvil, se plantea la construcción de una pista de pruebas. En la Figura \ref{fig:PPruebas}, se muestra la propuesta de pista de pruebas de un carril, plana y sin inclinación, en la que se pueden incorporar las señales de tránsito de vuelta a la derecha, vuelta a la izquierda, alto y un semáforo. La pista se plantea reconfigurable, con la posibilidad de cambiar la colocación de los señalamientos para validar la autonomía a través de distintas trayectorias. En esta ilustración, la línea roja representa una de las posibles trayectorias, mientras que los triángulos definen las localizaciones de las señales, mismas características que se pueden distribuir para tener trayectorias distintas. Las medidas finales de la pista están relacionadas con las dimensiones del vehículo, considerando, por ejemplo, el radio de giro.
\begin{figure}[htbp!]
	\centering
	\includegraphics[width=0.6\textwidth]{./Figuras/Pista/Pista}
	\caption{Pista de pruebas propuesta.}
	\label{fig:PPruebas}
\end{figure}
\par A priori, se plantea que esta plataforma móvil usará una configuración Ackerman, por el uso de la misma en vehículos de tamaño real. En el proyecto {\it AutoNOMOS Model} se cuenta con un chasís de un vehículo a control remoto que ha sido adaptado para portar los módulos sensoriales y de cómputo; por su efectividad comprobada, el empleo de algún chasís comercial es opción en la implementación, así como el partir de un chasís de diseño {\it Open-source}. Independientemente de la solución encontrada, se realizarán los análisis y simulaciones que permitirán aprovechar las cualidades intrínsecas de la plataforma, permitiendo además obtener parámetros de funcionamiento que servirán al momento de implementar el control del sistema.
\par Debido a que el proyecto se centra en el diseño e implementación de algoritmos para la conducción autónoma en escenarios urbanos a escala, por la amplitud del proyecto no se trabajará en la descripción de hardware ni la programación a bajo nivel, con lo cual se descarta la posibilidad de usar FPGA\footnote{Por sus siglas en inglés Field-Programable Gate Array} o DSP\footnote{Por sus siglas en inglés Digital Signal Processor}; con lo cual, una arquitectura que resulta cómoda por su compatibilidad con sistemas operativos Linux, es la ARM con ayuda de las computadoras de placa reducida de bajo costo, como las que han emergido en los años recientes, tomando con fines ilustrativos a la Raspberry Pi \cite{foundationRaspberryPiTeach}, Orange Pi \cite{OrangePiOrangepi}, Banana Pi \cite{BananaPiBPI} u Odroid \cite{ODROID}.\\
Para la realización de este proyecto, en la parte de software, se propone usar las librerías OpenCV \cite{opencvOpenCVReferenceManual} y OpenNI \cite{openniOpenNIProgrammerGuide} con el lenguaje de programación C++, por la compatibilidad existente con los sistemas operativos Linux sobre las arquitecturas mencionadas con anterioridad, y con diversos sensores RGB-D del mercado, como el Kinect \cite{InformacionSensorKinect}, el Asus Xtion \cite{XtionPROLIVE}, el Intel RealSense \cite{TecnologiaIntelRealSense}, o el Orbbec Astra \cite{gordonAstraSeries}. 
	\chapter{Desarrollo del trabajo}
	\label{cap:des}
	Para este Capítulo, tomando como base la metodología de trabajo propuesta en el protocolo de investigación, se describirá el proceso seguido para la realización de este Trabajo Terminal. Se inicia con una continuación del planteamiento del problema visto en el \autoref{cap:pdp}, en el que se realiza un análisis de los requerimientos previos al desarrollo del problema, en donde se exponen los requisitos del sistema y los subsistemas que componen el vehículo; esto se lleva a cabo en la Sección \ref{sec:prev}. En la Sección \ref{sec:mc} se extiende el modelo cinemático que se presentó como parte del marco teórico en la Sección \ref{sec:mc}, y se complementa el trabajo realizado con una serie de simulaciones que permitieron validar el control elegido para la plataforma como parte de la conducción autónoma. En las secciones que van desde la \ref{sec:selplat} hasta la \ref{sec:platpru}, se aborda la selección de los componentes que se emplearán en la plataforma robótica; tomando en consideración primero el chasís y elementos mecánicos, después los componentes de medición y procesamiento, pasando por los componentes electrónicos y finalizar en los elementos de la plataforma de pruebas; además de la adquisición de estas partes. Al final en la Sección \ref{sec:alg} se dispondrá de la documentación, a través de pseudocódigo, máquinas de estado y diagramas de flujo, de los algoritmos de visión artificial y control que se desarrollará en Trabajo Terminal. Quedando además la Sección \ref{sec:monveh} para documentar el montaje del vehículo.
\section{Requerimientos previos}
\label{sec:prev}
Desde el Protocolo de Investigación asociado a este proyecto se propusieron los subsistemas que se consideraron necesarios para el vehículo, teniendo junto con esto los principales requerimientos de cada uno de ellos. Se realizaron modificaciones sustanciales a estos requerimientos para adaptarlos a la perspectiva que se formuló para el robot. A pesar de no estar contemplada como una actividad en el cronograma de actividades, se realizó una revisión de este par de componentes metodológicos antes de iniciar de manera formal con las acciones propuestas en los cronogramas de actividades.
\subsection{Requerimientos del sistema}
\label{ssec:req}
Para desempeñar los objetivos planteados, el vehículo autónomo a escala tiene que cumplir con una serie de requerimientos. A continuación se resaltan con letras itálicas los requerimientos encontrados, mientras que la solución propuesta se presenta en letra normal.
\begin{itemize}
	\item {\it Es necesario capturar imágenes desde la cámara y enviarlas en tiempo real a la computadora a bordo}. El vehículo tendrá una cámara montada en la parte delantera sobre el chasís, y se comunicará de manera directa con la computadora a bordo.
	\item {\it La computadora necesita ser capaz de transmitir las órdenes a los actuadores}. Se puede contar con la computadora como un cerebro maestro que delegue enviar las señales a los actuadores a través de un sistema esclavo basado en microcontrolador.
	\item {\it Se ocupa contar con un chasís sobre el cual se monten los componentes del vehículo.} Por el alcance del trabajo, se dificulta realizar el diseño de un chasís, por lo que se adquirirá un vehículo a control remoto comercial como chasís, dicho vehículo deberá contar con los actuadores ya instalados; a este vehículo se le unirán el resto de los componentes de electrónicos que conforman la plataforma.
\end{itemize}
\subsection{Subsistemas del vehículo}
\label{ssec:subs}
El vehículo autónomo contará con cuatro subsistemas, con lo cual se puede delinear de manera particular su funcionamiento que si se representara como un único sistema. Lo que se muestra a continuación es la descripción de cada subsistema a través de sus tareas, lo que durante la programación ayudará al ser definidos como objetos.
\begin{itemize}
	\item {\bf Subsistema Cámara.} Se integra por el sensor RGB-D y se encarga de brindar mediciones como realimentación visual al subsistema de Control. El procesamiento se realiza en la computadora, y contiene las funciones para la identificación del carril y el reconocimiento de señalamientos de tránsito y semáforos.
	\item {\bf Subsistema Carro.} Se compone de la plataforma robótica como un conjunto de elementos mecánicos que trabajan en conjunto para realizar los movimientos especificados. Trabajan en conjunto los motores (servomotor de dirección y motor de desplazamiento). La computadora envía mediante comunicación serial los comandos de movimiento a un microcontrolador, que a su vez asigna esas configuraciones a los motores.
	\item {\bf Subsistema Control.} Coordina la ejecución de los movimientos enviando la información a los actuadores. Funciona compensando la medición del error obtendia en el subsitema de visión. El procesamiento se lleva a cabo en la computadora.
	\item {\bf Subsistema Navegación.} Se apoya de la odometría visual como herramienta para trazar el recorrido realizado por el robot. La odometría visual a su vez se auxilia de la IMU para validar sus propios datos.
\end{itemize}
\section{Modelo cinemático}
\label{sec:mc}
La principal característica (y quizá limitante) de los vehículos con la configuración Ackerman (configuración empleada por el carro común) es que no pueden desplazarse libremente en el plano, ya que no pueden rotar sobre su propio eje o desplazarse lateralmente; de manera exclusiva pueden avanzar o retroceder con cierta dirección. A esta clase de vehículos se les denomina no-holónomos. La configuración de los elementos mecánicos definen el modelo cinemático. Más adelante se verá que el modeo del sistema incluye también el modelo de los actuadores; el vehículo utiliza motores de corriente directa para actuar sobre los elementos mecánicos de la tracción y la dirección. Tanto el modelado de la plataforma Ackerman como el controlador propuesto se apoya en lo expuesto por G. Oriolo et al. en \cite{laumondRobotMotionPlanning1998}.
\par En primer lugar, se realizan las siguientes afirmaciones acerca de las condiciones sobre las que operará el robot, esto con el fin de permitir desarrollar la cinemática del robot móvil:
\begin{enumerate}
	\item El robot se encuentra sobre una superficie plana.
	\item Los ejes de guiado son perpendiculares al suelo.
	\item No existe deslizamiento entre la rueda y el suelo.
	\item El robot es un sólido rígido.
\end{enumerate}
\par Además de las afirmaciones anteriores, el modelo cinemático del vehículo se puede simplificar si se considera que la geometría de la configuración Ackerman permite llevar el modelo del vehículo al modelo de una bicicleta, esto considerando que el RMR está compuesto por un par de ruedas fijas en el eje trasero y un par de ruedas orientables en el eje delantero, suponiendo que las ruedas de cada eje se juntan en una sola en el punto medio de cada eje y su orientación es la orientación promedio de las ruedas presentes en el eje.
\par Se cuenta con un sistema fijo $(X, Y)$, que es el plano sobre el cual se encuentra el espacio de configuraciones del robot móvil. La configuración del vehículo está representada por sus coordenadas generalizadas $\mathbf{q}=[x, y, \theta, \varphi]^{T}$, en donde el par $(x, y)$ representa la posición del centro del eje trasero del robot, y estas coordenadas se miden con respecto al sistema de referencia fijo $(X, Y)$, el ángulo $\theta$ se forma entre el plano sagital del robot y el eje $X$, mientras que el ángulo $\varphi$ es la dirección promedio de las ruedas delanteras con respecto al plano sagital del robot. 
\par Un diagrama adecuado para representar cualquier configuración alcanzable del robot sobre el plano es el que se observa en la Figura \ref{fig:pose}. La velocidad lineal del vehículo está formada por las proyecciones de $(\dot{x}, \dot{y})$ de $v$ sobre el sistema $(X, Y)$. La velocidad angular $\omega$ actual sobre el punto $(x, y)$ del vehículo.
\begin{figure}[htbp!]
	\centering
	\includegraphics[width=0.5\textwidth]{./Figuras/Carro/Carro}
	\caption{Representación de la posición de un robot móvil con configuración Ackerman.}
	\label{fig:pose}
\end{figure}
\par A saber que el par $(\dot{x}, \dot{y})$ se puede a la vez expresar en términos de $v$ y $\theta$ como $[v\cos(\theta), v\sen(\theta)]$, mientras que el modelo cinemático expresado de manera matricial es el señalado en la Ecuación \eqref{eq:mc}. 
\begin{equation}\label{eq:mc}
	\left[\begin{array}{c}
	\dot{x}\\
	\dot{y}\\
	\dot{\theta}\\
	\dot{\varphi}
	\end{array}\right]=
	\left[\begin{array}{c}
	\cos(\theta)\\
	\sen(\theta)\\
	\frac{\tan(\varphi)}{L}\\
	0
	\end{array}\right]v+
	\left[\begin{array}{c}
	0\\
	0\\
	0\\
	1
	\end{array}\right]\omega
\end{equation}
\par Del robot móvil, las variables a controlar serán su velocidad lineal y su velocidad angular con el vector $u=[v, \omega]^{T}$. En el sistema fijo, el movimiento del robot como bicicleta está sujeto a dos restricciones de movimiento, una para cada rueda, y son el par de ecuaciones que representa la condición de no deslizamiento, Ecuación \eqref{eq:ns}, y la condición de rodamiento puro, Ecuación \eqref{eq:pr}.
\begin{subequations}
	\begin{equation}\label{eq:ns}
		\dot{x}\sen(\theta+\varphi)-\dot{y}\cos(\theta+\varphi)-L\dot{\theta}\cos(\varphi)=0
	\end{equation}
	\begin{equation}\label{eq:pr}
		\dot{x}\sen(\theta)-\dot{y}\cos(\theta)=0
	\end{equation}
\end{subequations}
\subsection{Diseño del controlador}
\label{ssec:dcon}
\par Desde el punto de vista de un conductor humano, la conducción de un vehículo requiere contar con un volante para ajustar la dirección de las ruedas frontales y de un acelerador para regular la velocidad lineal del vehículo. Si se pretende sustituir al conductor humano por una computadora, es necesario entonces establecer un rumbo que permita al controlador manipular la dirección y el acelerador de la manera que lo haría un conductor humano, e incluso mejor. 
\par En la Sección \ref{sec:con} se abordaron los diferentes enfoques de control posibles que se consideraron para este vehículo. Por el trabajo que debe realizar el RMR, la navegación por medio de referencias visuales en el terreno representa una opción para llevar a cabo el control de la plataforma. La navegación de vehículos autónomos se realiza en muchos casos controlando la posición respecto a una trayectoria generada con anterioridad; sin embargo para realizar la planeación de trayectorias y su posterior seguimiento es necesario contar con información previa del entorno sobre el que se navega, y es esta la razón por la que los vehículos autónomos incorporan un sistema GPS para realizar la planeación de la trayectoria con base en el terreno, y de igual manera en los vehículos autónomos a escala que se consultaron en el \autoref{cap:eda} se emula el sistema GPS con una cámara que apunta a un mapa ubicado en el techo del recinto donde se encuentra la pista. Por los alcances propuestos para el proyecto no se cuenta con sistema GPS ni su equivalente para vehículos a escala, con lo cual resulta de mayor conveniencia realizar la navegación reactiva basada en la identificación de elementos en el ambiente; lo que implica mantener al vehículo en el carril por el que se desplaza, además de realizar la identificación de los señalamientos de tránsito.
\par Continuando con la alusión al \autoref{cap:eda}, los proyectos que encabezan las investigaciones de la conducción autónoma suelen emplear metodologías basadas en controladores no lineales, tema que no se abordará en este trabajo. En el flujo de las tareas del vehículo, una vez que las imágenes son obtenidas por la cámara, la identificación de las líneas que limitan al carril es el primer paso del control, a partir de lo cual el error de orientación es la realimentación visual dada al sistema de control. Por un lado, una aproximación emplea una red neuronal que pueda decidir sobre la siguiente acción del vehículo, orientando las ruedas y regulando la velocidad; desde otra perspectiva, un controlador basado en lógica difusa permitiría orientar el vehículo a partir de las funciones de pertenencia que intervienen a la entrada del sistema; en la bibliografía se pueden encontrar casos de control lateral basados en los dos enfoques anteriores, sin embargo las arquitecturas de controladores basadas en PID resultan de mayor revelancia por su integración con el {\it visual servoing} que se encuentra en la bibliografía consultada.
\par Se ha tomado un controlador basado en PID para el seguimiento y permanecimiento en el carril. Con este fin es utilizado el modelo cinemático de la bicicleta para simplificar el problema, y a partir de él se definen las variables de error involucradas en el proceso de permanecimiento en el carril. Considerando una trayectoria establecida en el terreno con una curvatura suave, el vehículo deberá realizar un recorrido paralelo a la trayectoria y lo más cercano a ella. Se cuenta con la Figura \ref{fig:lane} en la que se describen las variables involucradas en el {\it controlador lateral} del vehículo, resaltando a $e$ como el error lateral y a $\delta$ como el error de orientación. El objetivo del controlador implica mantener al vehículo de manera paralela a la trayectoria logrando que el error de orientación sea mínimo, mientras que para lograr que el recorrido paralelo a la trayectoria sea cercano a ella se requiere que el error lateral sea también mínimo. Las variables $e$ y $\delta$ se pueden relacionar con las ecuaciones \eqref{eq:eibvs} y \eqref{eq:dibvs} para el caso del control IBVS o con las ecuaciones \eqref{eq:epbvs} y \eqref{eq:dpbvs} para el caso del control PBVS. De manera objetiva, el vehículo extrae características sólo de la imagen sin importar de manera específica la posición de la propia plataforma, por lo que el controlador IBVS resulta conveniente. El vehículo se encontrará viajando a lo largo del carril, y al ser $e$ de orden mínimo será suficiente controlar $\delta$ a través de $\varphi$, que es el viraje de las ruedas delanteras del vehículo.
\begin{figure}[htbp!]
	\centering
	\includegraphics[width=0.5\textwidth]{./Figuras/Lane}
	\caption{Generación del error para un controlador lateral de vehículo autónomo.}
	\label{fig:lane}
\end{figure}
\par Considerando en primer lugar una velocidad constante del vehículo cuando se desplaza sobre el eje $X$ mientras se estabiliza su orientación a través de oscilaciones con amplitud en el eje $Y$; se puede relacionar la distancia recorrida en el eje $X$ con el tiempo transcurrido durante esta acción. Se realizaron una serie de simulaciones aproximando la velocidad del vehículo en la simulación a la del vehículo real (de manera aproximada se ha establecido en *$5.7 \frac{cm}{s}$), tras lo cual se intentaron distintas ganancias $K_{p}$, $K_{d}$ y $K_{i}$ del controlador a modo de proceso de sintonización. Hay que aclarar que el vehículo real contará con valores diferentes en las ganancias en el controlador, ya que las simulaciones no contemplan las perturbaciones producidas por el terreno, como la fricción.
\begin{figure}[htbp!]
	\centering
	\includegraphics[width=0.7\textwidth]{./Figuras/PIDSim}
	\caption{Simulación del controlador PID.}
	\label{fig:pids}
\end{figure}
\par Se realizaron una serie de simulaciones para sintonizar el controlador de manera inicial. En la Figura \ref{fig:pids} se muestra una de estas simulaciones, en ella se muestra la configuración inicial del robot sobre el plano, además de la configuración final del robot al estabilizarse. En azul se muestra la trayectoria esperada, mientras que en naranja se muestra la trayectoria real. Las ganancias empleadas son *$K_{p}=2.7$, $K_{i}=0.001$ y $K_{d}=5.0$, y aunque el sistema presenta oscilaciones, son de un orden de unos pocos centímetros y no representan un gran problema para la cinemática del vehículo, se establecieron de esta manera para evitar que fueran de mayor magnitud. Los puntos a los que se establece que debe llegar el vehículo están relacionados con la curvatura local del carril, y son obtenidos gracias a la cámara, se toman parámtros importantes de las líneas como su pendiente e intercepto. La simulación que se presenta se realizó en el lenguaje de programación Python, y el listado del código de la misma puede encontrarse en el Anexo \ref{sec:csim}.
\subsubsection{Comportamiento eléctrico}
\label{sssec:elecomp}
\par Existen técnicas de control específicas para motores de CD, sin embargo --y para fortuna-- los carros RC comerciales en su mayoría incorporan un módulo ESC\footnote{Electronic Speed Controller.} que se encarga del control de velocidad del motor CD a través del ancho de pulso que recibe. Es de esta forma que es posible fijar distintos valores de velocidad con distintos anchos de pulso. Este es un funcionamiento similar al de los servomotores de posición. En la Figura \ref{fig:pwm} se muestran algunos de los anchos de pulso principales para un servomotor de posición, donde el estándar bajo el que trabajan involucra una frecuencia base de 50 Hz, logrando así que un servomotor con un pulso de 10 ms alcance una posición de $-90^{o}$, 15 ms equivale a $0^{o}$ y 20 ms equivale a $90^{o}$; en el caso de un motor de CD controlado por un ESC, 10 ms representa la máxima velocidad que alcanza el motor en reversa, 15 ms representa una velocidad cero y 20 ms representa la velocidad máxima que alcanza el motor hacia adelante.
\begin{figure}[htbp!]
	\centering
	\includegraphics[width=0.5\textwidth]{./Figuras/PWM}
	\caption{Representación de las diferentes señales enviadas a un servomotor.}
	\label{fig:pwm}
\end{figure}
\section{Selección de la plataforma}
\label{sec:selplat}
Dentro de los alcances y objetivos de este proyecto no se encuentran contemplados ni el diseño ni la construcción del chasís del robot, de esta manera es que se vuelve necesario realizar una búsqueda entre los múltiples chasís comerciales, para de esta forma tomar aquellos que se adapten a las necesidades de este Trabajo Terminal. Es habitual que el carro esté ligado a un control remoto que envíe señales para controlar las ruedas de acuerdo a los controles manipulados por el usuario; en este caso se desprecia la teleoperación y en su lugar se controla el vehículo a través de un programa dentro de la computadora embebida, que se encarga de ejecutar los módulos encargados de la interacción con los susbsistemas del vehículo.
\par Dentro de toda la comunidad del modelismo a escala, los vehículos a escala 1:10 son muy populares, y los hay para distintos objetivos. Se pueden encontrar vehículos especializados en velocidades altas y acrobacias sobre terrenos difíciles, los hay también especializados en derrapes y otros sólo de exhibición; además existen clasificaciones generales de acuerdo a su propulsión (existiendo eléctrica con y sin escobillas o de combustión interna).
\par Se tomaron en cuenta las afirmaciones expuestas en la Sección \ref{sec:mc} para eliminar de la selección aquellos vehículos que se encuentren fuera de los propósitos de este proyecto. Por incio de cuentas, como el vehículo se encuentra sobre una superficie plana, entonces no se toman en cuenta los que están diseñados para ir por terrenos accidentados; de acuerdo con la tercera afirmación, no existe deslizamiento entre la rueda y el suelo, por lo que los vehículos diseñados para realizar derrapes no se consideran. Fueron considerados sólo aquellos vehículos de propulsión eléctrica. Como el vehículo opera a velocidades bajas o moderadas sobre un terreno plano, se optó por elegir un vehículo de propósito general. En las tiendas en línea especializadas en vehículos RC se pueden encontrar a los mismos coches en un variado rango de precios, y como no se fija como requerimiento que tenga una alta resistencia a impactos, los chasís metálicos y de fibra de carbono quedan descalificados, y de paso ayuda con la reducción de costos. Es importante mencionar que la búsqueda se realizó dentro de una categoría llamada {\it ready to run} pues estos ya se encuentran preensamblados con un motor de tracción, servomotor de dirección y batería, que a la larga evita hacer múltiples selecciones de componentes.
\par La búsqueda descrita con anterioridad arroja como resultado preferir el modelo VRX-RH1025, cuya imagen se muestra en la Figura \ref{fig:chasis}, en el Anexo \ref{sec:vrx} toman parte los instructivos del ensamble mecánico del carro. Del chasís existen características significativas que a continuación se enumeran:
\begin{itemize}
	\item 364 mm de largo.
	\item 190 mm de ancho.
	\item 285 mm de distancia entre ejes.
	\item Transmisión de 6.5:1.
	\item Ruedas con radio de 65 mm y 26 mm de ancho.
	\item Batería incluida de 7.2 V y 1800 mAh.
\end{itemize}
\begin{figure}[htbp!]
	\centering
	\includegraphics[width=0.7\textwidth]{./Figuras/Chasis}
	\caption{Vista superior del chasís VRX RH1025}
	\label{fig:chasis}
\end{figure}
\par El vehículo cuenta con un módulo ESC modelo H0050B, con el fin de controlar el motor eléctrico con escobillas incorporado, del cual la hoja de datos se puede encontrar en el Anexo \ref{sec:ESC}.
\section{Selección de los componentes electrónicos}
\label{sec:selcomp}
La selección de los componentes electrónicos significa obtener requerimientos de diseño para la Sección \ref{sec:disel}. Los componentes que se seleccionan en esta etapa son la computadora del sistema, el sensor RGB-D y la IMU, al ser los elementos primordiales del sistema, con lo que es posible obtener su consumo eléctrico para la selección de la batería.
\subsection{IMU}
\label{ssec:velor}
Con base en proyectos similares al que aquí se desarrolla, para conocer la velocidad y orientación de vehículo se empleará un módulo de acelerómetro y giroscopio. Existen diferentes módulos que cumplen con esta función, sin embargo aquel que resulta ser de más fácil adquisición es el MPU6050 (Figura \ref{fig:MPU6050}), que involucra seis grados de libertad a través de sus tres ejes para el acelerómetro y otros tres ejes para la rotación con un giroscopio. Este módulo cuenta con una interfaz de comunicación I$^{2}$C, con la cual se transmitirań los datos hacia la computadora.
\begin{figure}[htbp!]
	\centering
	\includegraphics[width=0.35\textwidth]{./Figuras/MPU6050}
	\caption{Módulo de acelerómetro y giroscopio MPU6050}
	\label{fig:MPU6050}
\end{figure}
\par Gracias al uso de este sensor se vuelve posible resolver las restricciones en cuanto a las mediciones obtenidas por el sistema RGB-D. Una brújula provee una orientación absoluta respecto al polo norte magnético, mientras que un giroscopio provee el cambio en la orientación relativa con respecto al cambio en la orientación anterior. Tomando en consideración la sucesión de orientaciones con respecto a una fija e integrando en el tiempo, se puede contar con una orientación absoluta del carro.
\subsection{Sensor RGB-D}
\label{ssec:rgbd}
En realidad son pocos sensores RGB-D en el mercado, y la lista se limita aún más si se consideran sólo los que son compatibles con la librería OpenNI, por lo que es fácil incluir los principales en una tabla para mostrar de manera general sus características y a partir de ello realizar la selección del más adecuado. En la Tabla \ref{tab:rgbd} se consideran los sensores RGB-D principales de cada fabricante, pues dentro de un mismo fabricante se cuentan con características similares, quedan considerados los de características que atiendad a los requerimientos del proyecto.
\begin{table}[htbp!]
	\caption{Sensores RGB-D con sus principales características}
	\label{tab:rgbd}
	\begin{center}
		\begin{tabular}{|p{5cm}|p{7cm}|}
			\hline
			{\bf Cámara} & {\bf Características}\\
			\hline
			Orbbec Astra Pro \includegraphics[width=3.5cm]{./Figuras/Orbbec} &
			\begin{compactitem}
				\item 1280x720 pixeles a 30 FPS para RGB.
				\item 640x480 pixeles a 30 FPS para la profundidad.
				\item Rango de operación entre 0.6m y 8m.
				\item FOV de 60$^{O}$ en horizontal, 49.5$^{O}$ en vertical y 73$^{O}$ en diagonal.
			\end{compactitem}\\
			\hline
			Asus Xtion Pro \includegraphics[width=3.5cm]{./Figuras/Asus} &
			\begin{compactitem}
				\item 1280x1024 pixeles a 30 FPS para RGB.
				\item 640x480 pixeles a 30 FPS para la profundidad.
				\item Rango de operación entre 0.8m y 3.5m.
				\item FOV de 58$^{O}$ en horizontal, 45$^{O}$ en vertical y 70$^{O}$ en diagonal.
			\end{compactitem}\\
			\hline
			Intel Realsense SR300 \includegraphics[width=3.5cm]{./Figuras/Intel} &
			\begin{compactitem}
				\item 1920x1080 pixeles a 30 FPS para RGB.
				\item 1280x720 pixeles a 30 FPS para la profundidad.
				\item Rango de operación entre 0.2m y 1.5m.
				\item FOV de 69.4$^{O}$ en horizontal, 42.5$^{O}$ en vertical y 77$^{O}$ en diagonal.
			\end{compactitem}\\
			\hline
		\end{tabular}
	\end{center}
\end{table} 
\par La Tabla \ref{tab:rgbd} muestra los sensores RGB-D de los fabricantes Orbbec, Asus e Intel; para este proyecto no se consideró el Kinect de Microsoft debido a que el mismo se encuentra limitado en sus prestaciones de software y hardware al no contar con conexión USB. Para la selección de la cámara, la característica principal en la que se prestó atención es el rango de operación. El rango de operación determina la distancia mínima a la cual el sensor de profundidad garantiza encontrar puntos de interés, por lo que se requiere contar con un rango en el que el límite inferior sea cercano a la cámara, con lo cual el vehículo podrá operar en una pista también reducida. El límite superior del rango de operación no representa en este caso un parámetro determinante, ya que se trabajará con objetos próximos.
\par Con base en la búsqueda de la cámara con el límite inferior del rango de operación más reducido, se concluyó que las cámaras de la familia {\it Intel Realsense} son las que brindan un rango de operación inferior menor, con 0.2 m. El modelo SR300 de esa familia es el modelo más económico, y de mayor disponibilidad en los tiempos utilizables para este proyecto. El resto de las características técnicas son similares entre los diferentes fabricantes, además de que, aunque la cámara opere con una resolución de 1920x1080 pixeles o 1280x1024 pixeles, se reducirá la resolución de las imágenes para facilitar a la computadora realizar el procesamiento de los datos obtenidos, y el número de imágenes capturadas cada segundo se reducirá de acuerdo a la velocidad con la que navegue el RMR. En este mismo sentido, el FOV\footnote{Siglas en inglés para Campo de Visión} no difiere demasiado entre los tres modelos, y este parámetro no afecta en gran medida la operación del robot.
\par Una ventaja que se obtiene al seleccionar la cámara {\it Intel Realsense SR300} es que el propio fabricante Intel proporciona una SDK\footnote{Software Development Kit} para el manejo de las funciones primarias de la cámara, y ahora es posible la integración con la librería OpenCV. La cámara OpenCV requiere un puerto USB 3.0 para su operación, lo que se añade a los datos que se tomarán en cuenta para la selección de la computadora.
\subsection{Computadora de Placa Reducida}
\label{ssec:comp}
La labor de la computadora a bordo es procesar el volumen de datos que ha sido obtenido por el sensor RGB-D y la IMU. Basándose en los datos que aportan los sensores presentes es que se vuelve posible controlar la velocidad y orientación del vehículo. A partir de la versión 3.4 de la librería OpenCV es necesario contar con 1 Gb de memoria RAM para su instalación cuando se compila desde el código fuente (como en este caso que es necesario enlazar la librería OpenCV con las librerías LibRealSense y OpenNI); y en realidad para la operación posterior esta capacidad mínima depende de las operaciones a realizar, debido a que existen algoritmos bastante rápidos y otros que requieren emplear un mayor poder computacional. Con esto en mente se pre-seleccionaron las siguientes opciones debido a que cumplen con este requisito mínimo.
\par La Raspberry Pi en el modelo 3B cuenta con 1 Gb de memoria RAM, lo cual cumple con este requisito. Buscando contar con una mayor capacidad, se encontró en internet a las computadoras Odroid-XU4 y la Rock64, que representan fuertes candidaturas como computadora del vehículo. Se encontró que la computadora Odroid-XU4 no se encuentra disponible en la tienda oficial y no hay fecha para un relanzamiento; por otro lado en la tienda en línea de la computadora Rock64 se anuncia un nuevo lote disponible a la venta. La Tabla \ref{tab:computadoras} muestra las características de estas tres computadoras a manera de comparación.
\begin{table}[htbp!]
	\caption{Computadoras con sus principales características}
	\label{tab:computadoras}
	\begin{center}
		\begin{tabular}{|p{5cm}|p{7cm}|}
			\hline
			{\bf Computadora} & {\bf Características}\\
			\hline
			Raspberry Pi 3B \includegraphics[width=3.5cm]{./Figuras/RPi} &
			\begin{compactitem}
				\item Procesador ARM-8 con 1.2 GHz a 64 bits.
				\item 1 Gb de RAM.
				\item Funciona con memoria micro SD.
				\item Cuenta con conexión USB, Ethernet, WiFi, Bluetooth y GPIO.
				\item Consume 800mA a 5V
				\item Ya se cuenta con ella.
			\end{compactitem}\\
			\hline
			Odroid-XU4 \includegraphics[width=3.5cm]{./Figuras/Odroid} &
			\begin{compactitem}
				\item Procesador ARM-A7 con 2 GHz a 64 bits.
				\item 2 Gb de RAM.
				\item Funciona con memoria micro SD o eMMc.
				\item Cuenta con conexión USB, Ethernet y GPIO.
				\item Consume 4A a 5V.
			\end{compactitem}\\
			\hline
			Pine64-Rock64 \includegraphics[width=3.5cm]{./Figuras/Rock64} &
			\begin{compactitem}
				\item Procesador Rockchip RK3328.
				\item 4 Gb de RAM.
				\item Funciona con memoria micro SD o eMMc.
				\item Cuenta con conexión USB, Ethernet y GPIO.
				\item Consume 3A a 5V.
			\end{compactitem}\\
			\hline
		\end{tabular}
	\end{center}
\end{table}
\par En el Anexo \ref{sec:pi} se muestran los diagramas para los dos puerto GPIO que incorpora la tarjeta Rock64, que incluyen interfaces de entradas y salidas generales, además de comunicaciones UART, SPI e I$^{2}$C. Los otros periféricos disponibles y de utilidad son los siguientes:
\begin{itemize}
	\item Puerto Ethernet (1).
	\item Puerto USB 2.0 (2).
	\item Puerto USB 3.0 (1).
	\item Jack de audio 3.5 mm (1).
	\item Puerto HDMI (1).
	\item Puerto IR (1).
\end{itemize}
\subsection{Suministro de energía}
\label{ssec:ener}
Se cuenta con dos fuentes de energía separadas. El vehículo ya incluye una batería de 7.2 V a 1800 mAh, la cual le brinda al vehículo una autonomía de 25 minutos aproximadamente en el control de sus motores. Sin embargo para alimentar a la computadora (que a su vez alimentará vía USB a la cámara y vía GPIO a la IMU y a la tarjeta controladora de los motores) se requerirá una fuente de energía separada que brinde al menos la misma autonomía para la computadora.
\begin{figure}[htbp!]
	\centering
	\includegraphics[width=0.4\textwidth]{./Figuras/Power}
	\caption{Batería elegida para el sistema.}
	\label{fig:power}
\end{figure}
\par En la actualidad se cuenta con una amplia variedad de fuentes de energía portátiles, en especial resaltan aquellas empleadas como baterías externas para celulares y tabletas, que son también llamadas {\it powerbank}. Las baterías de las que se habla han adquirido altas prestaciones, con lo que pueden brindar hasta 3.1 A de corriente en algunos casos a través de un puerto USB.
\par Mediante una búsqueda realizada por internet se encontraron diferentes batería portátiles, de las cuales son útiles aquellas que puedan proveer 3.1 A y al menos los 25 minutos de autonomía, del carro. Con una capacidad de 3000 mAh se puede mantener la computadora encendida por una hora antes de recargar la batería, y esta característica resulta suficiente para una sesión de pruebas con el vehículo. La batería elegida suministra la corriente requerida y cuenta con una capacidad de 20000 mAh, suficiente para más de *seis horas de actividad de la computadora y la cámara.
\section{Diseño electrónico}
\label{sec:disel}
En esta Sección se detallan las decisiones de diseño que se tomaron a partir de los elementos más importantes de la selección de componentes: Chasís, sensor RGB-D, computadora, IMU y {\it powerbank}.
\subsection{Interfaz para el control de los motores}
\label{ssec:motcont}
Como ya se ha mencionado antes, para realizar el control de los motores será menester emplear dos señales de PWM; una para el ESC del motor de tracción y otra para el servomotor. Realizar esta acción desde la computadora se vuelve complicado, y esta es una razón suficiente para optar por el diseño de una tarjeta externa basada en microcontrolador como sistema esclavo, que facilite el objetivo de enviar ambas señales de PWM; se propone codificar las características de estas señales PWM para ser enviadas por la computadora mediante comunicación UART.
\par El microcontrolador ATMEGA328P se vende a un precio bajo en comparación de sus prestaciones (de las que sólo por enumerar se cuenta con funcionalidad para comunicación UART, SPI e I$^{2}$C, entradas y salidas generales, ADC, comparadores analógicos, timers de 8 y 16 bits y puertos con salida PWM), lo que lo hace viable para ser el cerebro de la tarjeta que se encargará de controlar los motores. El exceso de puertos de entrada y salida permitirá la escalabilidad del proyecto en un futuro.
\par El diseño de la PCB se realizó en el software Eagle de Autodesk. Lo que se verá a continuación son el esquemático, la vista de las pistas y la vista tridimensional de la PCB en las Figuras \ref{fig:Esque}, \ref{fig:BRD} y \ref{fig:PCB}. El circuito cuenta con un cristal externo para estabilizar la señal de reloj, además el programa que se grabó en el microcontrolador y que controla los motores es el mostrado en el Anexo \ref{sec:micro}.
\begin{figure}[htbp!]
	\centering
	\begin{subfigure}[htbp!]{0.4\textwidth}
		\includegraphics[width=\textwidth]{./Figuras/Esque}
		\caption{Circuito esquemático con las conexiones de la tarjeta.}
		\label{fig:Esque}
	\end{subfigure}
	\begin{subfigure}[htbp!]{0.25\textwidth}
		\includegraphics[width=\textwidth]{./Figuras/BRD}
		\caption{Pistas y conexiones del circuito impreso.}
		\label{fig:BRD}
	\end{subfigure}
	\begin{subfigure}[htbp!]{0.3\textwidth}
		\includegraphics[width=\textwidth]{./Figuras/PCB}
		\caption{Vista tridimensional de la PCB.}
		\label{fig:PCB}
	\end{subfigure}
	\caption{PCB para el control de los motores.}
\end{figure}
\subsection{Conexiones internas}
\label{ssec:conint}
El vehículo funciona de manera inalámbrica sin intervención del operador, sin embargo para este trabajo se implementa una interfaz de comunicación SSH entre la computadora Rock64 y la computadora del usuario para apreciar el procesamiento realizado por la computadora, esto sólo con fines ilustrativos. Para el vehículo es necesario realizar la interconexión de los diferentes módulos que conforman el sistema. El nodo central del sistema es la computadora Rock64, la cual se conecta por medio de un puerto USB 3.0 a la cámara Intel Realsense SR300, y por UART envía de manera codificada los anchos de pulso para el controlador de los motores, que maneja por PWM el ángulo del servomotor de dirección, mientras que también por PWM se maneja la velocidad y la dirección a la que debe girar el motor de tracción. La comunicación con la IMU se da a a través del microcontrolador mediante I$^{2}$C y se comunica la información a la computadora mediante UART. Se creó el diagrama de la Figura \ref{fig:conexiones} para ilustrar las manera en que interactuan las partes del sistema.
\begin{figure}[htbp!]
	\centering
	\includegraphics[width=0.5\textwidth]{./Figuras/Conexiones}
	\caption{Diagrama de conexiones internas del vehículo.}
	\label{fig:conexiones}
\end{figure}
\section{Diseño y construcción de la pista de pruebas}
\label{sec:platpru}
El vehículo seleccionado tiene 364 mm de largo y 190 mm de ancho. Otro dato que es importante de conocer son los límites de giro de las ruedas, estando ese rango en el vehículo entre -45$^o$ y 45$^o$. Con base en estos datos se pueden establecer los parámetros con los que debe cumplir la pista de pruebas. Las láminas de MDF de 3 mm de ancho tienen dimensiones de 1.22x2.44 m, con lo que las dimensiones finales de la pista serán múltiplos de 1.22 m.
\par Ya se ha mencionado desde la Sección \ref{ssec:mc} que el radio de curvatura para un vehículo Ackerman puede ser calculado por la expresión $R=L/\tan(\varPsi)$; se ha determinado de manera experimental que el ángulo de giro del vehículo se encuentra entre -45$^{o}$ y 45$^{o}$, de modo que al sustituir la distancia entre ejes y el ángulo de giro máximo se obtiene que el radio de giro para este vehículo en específico es $R=285mm/\tan(45^o)=285mm$, lo que indica que para dar una vuelta perpendicular el vehículo debe empezar a virar antes de que su borde frontal se encuentre en el borde más cercano del camino, de lo contrario se puede salir de la pista.
\begin{figure}[htbp!]
	\centering
	\includegraphics[width=0.8\textwidth]{./Figuras/Pista}
	\caption{Pista de pruebas propuesta acotada.}
	\label{fig:Pista}
\end{figure}
\par Se estableció desde el Protocolo de Investigación que la pista de pruebas contaría sólo con un carril, esto debido a que sólo se proponen maniobras individuales. Para la pista se considera el ancho del carril como 1.5 veces el ancho del vehículo, en este caso da como resultado 275 mm. En proyectos similares se emplean pistas de al menos 8 veces el largo del vehículo para realizar las pruebas. Con lo que se obtiene como resultado 2912 mm, y este número se redondeó para que el resultado encaje con las medidas de las láminas de MDF. Ahora, considerando el boceto de la pista que se propuso en la Figura \ref{fig:PPruebas}, se propone el siguiente croquis de la pista de pruebas. La pista mostrada en el boceto permite la circulación continua del vehículo, una clara ventaja con respecto a lo propuesto en el protocolo de investigación.
\begin{figure}[htbp!]
	\centering
	\includegraphics[width=0.8\textwidth]{./Figuras/PistaC}
	\caption{Pista de pruebas construida.}
	\label{fig:pista}
\end{figure}
\par Al inicio de las actividades realizadas para Trabajo Terminal II se encontró la construcción de la pista de pruebas, quedando como resultado la Figura \ref{fig:pista}, que es la unión de los tres tramos de pista.
\begin{figure}[htbp!]
	\centering
	\includegraphics[width=0.5\textwidth]{./Figuras/Senales}
	\caption{Señalamientos fabricados para la pista de pruebas.}
	\label{fig:Sen}
\end{figure}
\par Para la navegación a través del camino por medio de señalamientos el vehículo deberá reconocer las señales de tránsito que se fabricaron en, y que aparecen en la Figura \ref{fig:Sen}. Los señalamientos tienen una altura de 12 cm y están fabricados mediante impresión 3D con PLA. La imagen de la señalética está impresa en papel adherible.
\section{Montaje del vehículo}
\label{sec:monveh}
A continuación se hará una recapitulación de los componentes del vehículo.
\begin{itemize}
	\item {\bf Sensor RGB-D.} Para cumplir con este propósito se seleccionó la cámara Intel Realsense SR300 (Figura \ref{fig:SR300}), la cual cuenta con una resolución de 1920x1080 pixeles en el espacio de color RGB, presenta junto con eso una resolución de 1280x720 pixeles para la profundidad. El rango de operación de la cámara es entre 0.2 y 1.5 m con un campo de visión de 69.4$^o$ en horizontal, 42.5$^o$ en vertical y 77$^o$ en diagonal. A la cámara se le adaptará una zapata que la mantendrá fija al chasís del carro.
	\begin{figure}[htbp!]
		\centering
		\includegraphics[width=0.4\textwidth]{./Figuras/Intel}
		\caption{Sensor RGB-D Intel Realsense SR300}
		\label{fig:SR300}
	\end{figure}
	\item {\bf Computadora.} La plataforma robótica empleará una computadora Rock64 (Figura \ref{fig:Rock64}) de la compañía Pine64. La computadora Rock64 es una computadora de placa reducida con arquitectura ARM de 64 bits, emplea 4 Gb de memoria RAM, e incorpora un puerto para conectar una memoria microSD y/o eMMC para el almacenamiento. Como periféricos externos, la computadora incorpora dos puertos USB 2.0 y uno 3.0, además cuenta con conector ethernet, jack de audio de 3.5 mm, puerto HDMI, y un conector GPIO Pi-2 y Pi-5. Un elemento que será importante durante las pruebas e implementación de la navegación es el receptor IR incorporado. La computadora se alimenta a 5 V y 3 A mediante un conector de barril de 3.5 mm. El principal sistema operativo soportado por la computadora es el Ubuntu 18.04, sobre el cual se instalará la librería OpenCV y el Realsense SDK de Intel, que serán las encargadas de tomar los datos procedentes de la cámara y el posterior procesamiento.
	\begin{figure}[htbp!]
		\centering
		\includegraphics[width=0.4\textwidth]{./Figuras/Rock64}
		\caption{Computador Rock64 de la compañía Pine64}
		\label{fig:Rock64}
	\end{figure} 
	\item {\bf Chasís.} El chasís consiste en dos partes. La parte inferior es el chasís de un vehículo a control remoto a escala 1:10. Para este proyecto se optó por el vehículo VRX RH1025. El chasís superior es una base de acrílico de 3 mm cortada con láser la cual irán montados los demás componentes. El chasís inferior no es un juguete ordinario, pues es una representación realista de un carro real. Cuenta con cuatro ruedas y tracción trasera con una velocidad máxima de 15 km/h. El vehículo cuenta con un motor brushless DC para lograr su desplazamiento, e incluye un circuito ESC para controlarlo usando PWM; además el vehículo cuenta con un servomotor para el control de las ruedas en la estructura Ackerman; y junto con esto una batería de 7.2 V y 1800 mAh, proveyendo energía a estos sistemas. Una imagen del chasís inferior se puede encontrar en la Figura \ref{fig:chasis}; los componentes montados sobre el chasís superior se aprecian en la Figura \ref{fig:chas}, mientras que el plano de la placa puede apreciarse en el Anexo \ref{sec:chass}.
	\begin{figure}[htbp!]
		\centering
		\includegraphics[width=0.8\textwidth]{./Figuras/PlacaEns}
		\caption{Vista superior del chasís superior del vehículo}
		\label{fig:chas}
	\end{figure}
	\item {\bf Acelerómetro y giroscopio.} El módulo de acelerómetro y giroscopio que se tomó en consideración es el MPU6050, el cual cuenta con seis grados de libertad: Tres lineales y tres angulares. El módulo se comunica mediante el protocolo I$^2$C para la transmisión de datos. El módulo de acelerómetro y giroscopio se puede encontrar en la Figura \ref{fig:MPU6050}
	\item {\bf Interfaz para el control de los motores.} Para este propósito se diseñó una tarjeta de circuito impreso con el microcontrolador ATMEGA328P. El objetivo es proveer los puertos para manejar el PWM para el ESC del motor brushless y el servomotor. Los valores necesarios para el manejo de estos anchos de pulso será recibido por el microcontrolador desde la computadora Rock64. Con anterioridad se cuentan cn tres imágenes que muestran más detalle a la PCB, y son las Figuras \ref{fig:Esque}, \ref{fig:BRD} y \ref{fig:PCB}.
	\item {\bf Suministro de energía.} Se mencionó antes que el vehículo a control remoto que fue seleccionado incluye ya una batería que será de uso exclusivo para los motores. Con el objetivo de proveer energía a los sistemas como la computadora y los sensores, se incluye además una batería de que pueda suministrar hasta 3.1 A y 20000 mAh para brindar autonomía durante las pruebas que se realizarán.
\end{itemize}
\par Junto con esto se presenta una tabla con los costos que ha alcanzado el proyecto hasta este momento.
\begin{table}[htbp!]
	\caption{Costos condensados del proyecto.}
	\label{tab:costos}
	\begin{center}
		\begin{tabular}{|c|c|c|c|}
			\hline
			{\bf Componente} & {\bf Modelo} & {\bf Costo} & {\bf Moneda}\\
			\hline
			\hline
			Computadora & Rock64 & \$ 899 & MXN\\
			\hline
			Chasís & VRX RH1025 & \$ 2280 & MXN\\
			\hline
			Cámara & Intel Realsense SR300 & \$ 1780 & MXN\\
			\hline
			Acelerómetro y giroscopio & MPU6050 & \$ 100 & MXN\\
			\hline
			PCB & Control de motores & \$ 100 & MXN\\
			\hline
			Energía & {\it Powerbank} & \$ 399 & MXN\\
			\hline
			Pista & Pista de pruebas & \$ 400 & MXN\\
			\hline
			Extras & Chasís superior y base para cámara & \$ 150 & MXN\\
			\hline
		\end{tabular}
	\end{center}
\end{table}
\par Un proceso de importancia en el proceso fue la puesta a punto de la computadora y la cámara. El proceso involucra la instalación del sistema operativo en la tarjeta microSD, después instalar una larga lista de dependencias de software para culminar con la instalación de las librerías LibRealSense y OpenCV. Se encuentra con más detalle el procedimiento en el Anexo \ref{cap:pap}.
\par En cuanto a lo que concierne de manera general al montaje del vehículo, se inicia con el chasís inferior que está completamente armado de fábrica, y sobre él se coloca la placa que conforma al chasís superior. La placa cuenta con perforaciones para tornillo M3 para montar la computadora, la IMU y la placa controladora de motores. Las conexiones entre los diferentes componentes se realizan con cables Dupont, y se aseguran con cinchos. En las Figuras \ref{fig:fron}, \ref{fig:sup} e \ref{fig:iso} se observan vistas del vehículo ensamblado.
\begin{figure}[htbp!]
	\centering
	\begin{subfigure}[htbp!]{0.4\textwidth}
		\includegraphics[width=\textwidth]{./Figuras/CarroFron}
		\caption{Vista frontal.}
		\label{fig:fron}
	\end{subfigure}
	\begin{subfigure}[htbp!]{0.4\textwidth}
		\includegraphics[width=\textwidth]{./Figuras/CarroSup}
		\caption{Vista superior.}
		\label{fig:sup}
	\end{subfigure}
	\begin{subfigure}[htbp!]{0.4\textwidth}
		\includegraphics[width=\textwidth]{./Figuras/CarroIso}
		\caption{Vista isométrica.}
		\label{fig:iso}
	\end{subfigure}
	\caption{Diferentes vistas del vehículo armado.}
\end{figure}
\section{Algoritmos e implementación}
\label{sec:alg}
En esta Sección se discute la implementación de los algoritmos de conducción autónoma del vehículo, los cuales se desarrollan en la computadora a bordo. Para la navegación del vehículo se realizó una máquina de estados que en la que se abarca el proceso a partir de una condición inicial de arranque hasta una condición final de paro. Durante el trayecto existirán condiciones intermedias que obligarán a la computadora a tomar una decisión acerca de la acción más conveniente a realizar. Esta primera máquina de estados se muestra en la Figura \ref{fig:me1}.
\par En el diagrama mostrado, el estado {\bf Avanzar} es un nombre abreviado para el mantenimiento en el carril del vehículo, y a partir de ese estado emanan los siguientes. A los estados {\bf Derecha} e {\bf Izquierda} se accede al haber detectado el vehículo la correspondiente señal de tránsito, y vuelven al estado {\bf Avanzar} al momento en el que el vehículo ha concluido la maniobra. El estado {\bf Alto} es el que se obtiene cuando el vehículo se encuentra con un semáforo cuya luz encendida es la roja, y permanece en él hasta que la luz cambie a verde. Por último, el estado {\bf Final} sucede cuando el vehículo se encuentra con una señal de Alto.
\begin{figure}[htbp!]
	\centering
	\includegraphics[width=0.7\textwidth]{./Figuras/ME1}
	\caption{Máquina de estados para la navegación del vehículo.}
	\label{fig:me1}
\end{figure}
\par A continuación se desglosa cada una de las actividades en los estados y los propios algoritmos empleados para el mantenimiento en el carril y el reconocimiento de estructuras.
\subsection{Mantenimiento en el carril (Avanzar)}
\label{ssec:ava}
Esta acción se compone de varias etapas, siendo la primera el propio reconocimiento del carril. El reconocimiento del carril sigue la secuencia mostrada en la Figura \ref{fig:pipe}, donde el primer paso consiste en la propia captura de la imagen.
\begin{figure}[htbp!]
	\centering
	\includegraphics[width=\textwidth]{./Figuras/Pipe}
	\caption{Flujo de trabajo para el reconocimiento del carril}
	\label{fig:pipe}
\end{figure}
\par El proceso que desemboca en la obtención de las características de las líneas que definen los extremos del carril inicia con la imagen de color capturada por la cámara, a la cual se le aplica un filtro a la imagen capturada se le aplica un filtro gaussiano con un núcleo de dimensión (3x3) que reduce el ruido provocado por efectos de iluminación brillo por iluminación. La imagen filtrada pasa por un proceso de resaltamiento de bordes con el operador Canny, en el que la imagen resultante ahora sólo contiene elementos blancos y negros, de los cuales sólo es importante un segmento trapezoidal que contiene únicamente el carril. En la imagen de la región de interés ahora se pueden buscar líneas gracias a la transformada de Hough, sin embargo el número de líneas detectadas es alto, y se requiere  filtrar aquellas que son de importancia; un criterio para realizar la separación es de acuerdo a su pendiente, puesto que las líneas a la izquierda del carril tienen pendiente negativa y las que se encuentran a la derecha tienen pendiente positiva; a ambos lados se almacenan en un vector. Por último, de las características de las líneas a ambos lados se puede realizar una interpolación lineal para ajustar sólo una línea a cada lado del carril, de la que se conoce su pendiente y su intercepto; la imagen mostrada al final de la Figura \ref{fig:pipe} muestra las pendientes a cada lado del carril junto con el área encerrada entre ellas.
\par Como siguiente paso en la preservación del carril, se toman las pendienes de ambas líneas y se suman (una pendiente es positiva y la otra negativa), siendo el resultado la pendiente en la línea que une la intersección de las dos rectas con el centro del inicio de ambas líneas. El controlador toma como entrada esta pendiente y busca llevarla a cero a lo largo de todo el recorrido, como se pantea en la Sección \ref{ssec:dcon}.
\par En las simulaciones se presentaron las ganancias tomadas para el carro, sin embargo para el carro real fue necesario recalcular las ganancias, y el proceso se realizó en tres etapas: primero se hizo teniendo el vehículo estático mientras reconocía el carril, y se llevaron las ganancias del controlador PID hasta que en cada iteración del algoritmo había menos de *0.01 \% de error; la segunda etapa resultó similar, teniendo la diferencia de que ahora el servomotor encargado de la dirección intenta seguir la línea media del carril; la tercera etapa consistió en corregir las oscilaciones mientras que vehículo se encontraba avanzando en un tramo recto de pista a una velocidad constante. Las ganancias obtenidas son: *$K_{p}$, $K_{d}$ y $K_{i}$ 6.5, 10.0 y 0.0001, como corresponde.
\subsection{Reconocimiento de estructuras}
\label{ssec:rde}
Considerando que el sensor RGB-D provee una imagen como:
\begin{itemize}
	\item Una imagen a color en formato RGB.
	\item Un mapa de profundidad con las intensidades que representan distancia.
\end{itemize}
\par Se puede  tomar como referencia el trabajo publicado en \cite{herediafavieriReconocimientoObjetosImagenes2015}, en el que se cuenta con algoritmos para relacionar datos de 2D y de 3D. El sensor RGB-D proporciona un cuadro con más de 300, 000 puntos. Tal y como ocurre en la mayoría de las aplicaciones de la robótica móvil, se puede simplificar el proceso si el robot se considera plano y su forma viene determinada por el mayor polígono que circunscribe las distintas secciones en altura de éste. Una solución eficiente y que no compromete la eficacia del algoritmo reactivo  utilizado consiste en reducir la información 3D obtenida y proyectarla sobre el plano 2D en el que se encuentra el robot. La simplificación resultante consiste en un barrido horizontal que contiene las mediciones más cercanas a los obstáculos presentes a diferentes alturas.
\par Puede ser útil segmentar el mapa de profundidad a partir de la información tomada por la imagen a color. La detección de elementos como semáforos y señalamientos de tránsito se puede realizar mediante el emparejamiento de plantillas a partir de imágenes de prueba, con lo que los algoritmos de FAST, SIFT, SURF y otros, resultan útiles para encontrar la región de la imagen donde se encuentra el señalamiento, y encontrar la distancia a él gracias al mapa de color.
\par El subsecuente proceso para la conversión de datos involucra asumir que se cuenta ya con una imagen a color y un mapa de profundidad provistos por la cámara; el proceso se basa en el algoritmo SURF, a partir de lo cual se sigue el siguiente procedimiento de la manera siguiente:
\begin{itemize}
	\item Detectar características en la imagen objeto (este procedimiento se realiza para cada señal).
	\item Detectar características en la imagen objeto (este procedimiento se realiza para cada cuadro).
	\item Emparejar las características entre ambas imágenes y se calcula el valor de correlación.
	\item Si el nivel de correlación es mayor a 0.80, se obtiene el centro de la imagen, se calcula la profundidad en ese punto y se dibuja en la imagen el área que delimita a la señal en la escena. De lo contrario continua buscando correlaciones en la imagen.
\end{itemize}
\begin{figure}[htbp!]
	\centering
	\includegraphics[width=0.3\textwidth]{./Figuras/SURFDere}
	\includegraphics[width=0.3\textwidth]{./Figuras/SURFAlto}
	\includegraphics[width=0.3\textwidth]{./Figuras/SURFIzq}
	\label{fig:surf}
	\caption{Aplicación del algoritmo SURF para la detección de señales.}
\end{figure}
\par En las pruebas realizadas se comprobó la eficacia del algoritmo SURF para la detección de los señalamientos de tránsito, las pruebas se encuentran en la Figura 4.19, donde en la esquina superior izquierda se encuentra la imagen objeto, mientras que la imagen escena representa la mayor parte de la imagen. Las líneas de múltiples colores unen los descriptores de ambas imágenes, mientras que el cuadrilátero azul muestra la localización de la imagen objeto dentro de la escena.
\par Para la última parte del procedimiento, que involucra la distancia al objeto, se muestra con la Figura \ref{fig:dist} la distancia al centro de una señal de alto dentro de la imagen de profundidad.
\begin{figure}[htbp!]
	\centering
	\includegraphics[width=0.5\textwidth]{./Figuras/Distancia}
	\caption{Distancia al señalamiento.}
	\label{fig:dist}
\end{figure}
\par El procedimiento es similar para la detección de semáforos. Una vez que la imagen objeto (semáforo de prueba) se empareja en la imagen de escena, se segmenta la imagen por color de acuerdo al que se busca (Rojo, Verde o Amarillo) y se buscan círculos en la imagen. Un ejemplo del procesamiento se muestra en la Figura \ref{fig:sema}, resaltando la luz roja del semáforo. El dibujo del círculo para señalar la luz se realizó en la imagen segmentada en color con el objetivo de lograr que se aprecie.
\begin{figure}[htbp!]
	\centering
	\includegraphics[width=0.5\textwidth]{./Figuras/Sema}
	\caption{Localización de luz roja de un semáforo.}
	\label{fig:sema}
\end{figure}
\subsection{Programación}
Todos los algoritmos que se observan a lo largo de la sección se programaron en C++ usando la librería OpenCV (enlazada con LibRealSense y OpenNI). Se dividió el problema principal en subsistemas, y de cada uno de ellos se construyó una clase compuesta por un archivo de cabecera y uno de funciones. Se cuenta en la Figura \ref{fig:dep} con las dependencias entre cada uno de las componentes de software, teniendo el la parte superior al progama principal y a los archivos de prueba de cada uno de los módulos.
\begin{figure}[htbp!]
	\centering
	\includegraphics[width=0.5\textwidth]{./Figuras/Dependencias}
	\caption{Diagrama de dependencias de los componentes.}
	\label{fig:dep}
\end{figure}
\par En los Anexos \ref{sec:cam}, \ref{sec:cont}, \ref{sec:car} y \ref{sec:main} se presentan los subsistemas del sistema como clases y además se presenta el programa principal. Se omiten los archivos de prueba de los módulos. Los diagramas UML de las cuatro clases principales del proyecto se presentan en las Figuras \ref{fig:camC}, \ref{fig:controlC} y\ref{fig:carroC}.
\begin{figure}[htbp!]
	\centering
	\includegraphics[width=0.3\textwidth]{./Figuras/carroClass}
	\caption{Clase Carro}
	\label{fig:carroC}
\end{figure}
\par La Clase Carro involucra las variables obtenidas por la IMU, y también trabaja con el UART del microcontrolador.
\begin{itemize}
	\item {\it initUART}. Inicializa la comunicación UART entre la tarjeta con el microcontrolador.
	\item {\it cerrar}. Termina la comunicación UART.
	\item {\it setVal}. Envía al carro los valores de la velocidad y el ángulo de las ruedas.
\end{itemize}
\begin{figure}[htbp!]
	\centering
	\includegraphics[width=0.7\textwidth]{./Figuras/camClass}
	\caption{Clase Cámara.}
	\label{fig:camC}
\end{figure}
\par La clase maneja la información relacionada al sensor RGB-D.
\begin{itemize}
	\item {\it filtro}. Aplica filtro gaussiano a la imagen de entrada para remover el ruido 
	\item {\it blanco}. Convierte la imagen a HSV para detectar las líneas blancas de la imagen.. 
	\item {\it bordes}. Detecta todos los bordes en la imagen filtrada. 
	\item {\it RDI}. Aplica una máscara poligonal a la imagen para reducir la región de interés. 
	\item {\it detectarL}. Emplea la transformada de Hough probabilística para obtener todos los segmentos de línea en la región de interés.
	\item {\it sperararL}. Separar las líneas a la izquierda y a la derecha del carril 
	\item {\it MinCu}. Se encarga de tomar todas las líneas a ambos lados del carril para interpolar una sola línea a cada lado del carril. 
	\item {\it dibujarC}. Se encarga de mostrar las líneas detectadas para el carril junto con el área encerrada entre ellas. 
	\item {\it detectarDescO}. Se encarga de tomar las imágenes y por medio del algoritmo SURF encontrar los descriptores en la imagen de objeto. 
	\item {\it detectarDescE}. Encuentra los descriptores de la escena. 
	\item {\it match}. Empata los descriptores del objeto con los de la escena. 
	\item {\it dibujarMatch}. Muestra los matches entre el objeto y la escena de manera gráfica. 
	\item {\it locEscena}. Relaciona las transformaciones geométricas entre los descriptores del objeto con los descriptores de la escena. 
	\item {\it limitesO}. Determina los límites del objeto y su correspondencia en la escena. 
	\item {\it dibujarO}. Dibuja al objeto encerrado dentro de la escena. 
	\item {\it Alto}. Condensa los parámetros necesarios para buscar la señal de alto en la escena. 
	\item {\it Izquierda}. Condensa los parámetros necesarios para buscar la señal de vuelta a la izquierda en la escena. 
	\item {\it Derecha}. Condensa los parámetros necesarios para buscar la señal de vuelta a la derecha en la escena.
	\item {\it getDistancia}. Obtiene la distancia del punto p con la información de la imagen de profundidad.
	\item {\it dibujar}. Se encarga de mostrar las líneas detectadas para el carril junto con el área encerrada entre ellas. Muestra los señalamientos encerrados y muestra las distancias al final del camino detectado y a los señalamientos. 
\end{itemize}
\begin{figure}[htbp!]
	\centering
	\includegraphics[width=0.6\textwidth]{./Figuras/controlClass}
	\caption{Clase Control}
	\label{fig:controlC}
\end{figure}
\par La clase maneja la información relacionada al controlador PID. El constructor toma los valores deseados de las ganancias $K_{p}$, $K_{I}$ y $K_{D}$.
\begin{itemize}
	\item {\it init}. Inicializa las variables globales de la clase con los valores locales dados.
	\item {\it error}. Actualiza los valores del error en el controlador.
	\item {\it tErr}. Emplea todos los errores y las ganancias para calcular la compensación.
\end{itemize}
	\chapter{Análisis y validación de resultados}
	\label{cap:resul}
	El proceso de pruebas para este proyecto se realizó a través de distintas etapas. De manera estática, se realizó la detección de las líneas del carril. En el proceso se encontró que la imagen recibía ruido procedente del color del suelo adyacente a la pista, lo cual fue posible mitigar al controlar las condiciones en las periferias del escenario. Otro error recurrente durante la detección del carril se obtiene cuando la cámara no se encuentra en una posición recta, provocando que sólo se detecte una de las líneas del carril.
\par La detección de los señalamientos de tránsito no es un proceso robusto, al contar sólo con una imagen para emparejamiento para cada señal. El algoritmo SURF aumenta sus costos computacionales cuando el emparejamiento toma un número mayor de imágenes objeto, lo cual dificulta al vehículo identificar los señalamientos por el tiempo que toma a la computadora procesar la información. Como desventaja de emplear sólo una imagen objeto para cada señal, es que la si la inclinación del plano del señalamiento tiene una apertura amplia con respecto al plano del vehículo, entonces la adquisición y comparación de descriptores no es existosa y, por lo tanto, el vehículo no detecta el señalamiento. En el par de imágenes de la Figura \ref{fig:SSURF} se muestran emparejamientos correctos con el algoritmo.
\begin{figure}[htbp!]
	\centering
	\begin{subfigure}[htbp!]{0.4\textwidth}
		\includegraphics[width=\textwidth]{./Figuras/SSURF1}
	\end{subfigure}
	\begin{subfigure}[htbp!]{0.4\textwidth}
		\includegraphics[width=\textwidth]{./Figuras/SSURF2}
	\end{subfigure}
	\caption{Emparejamientos positivos con el algoritmo SURF.}
	\label{fig:SSURF}
\end{figure}
\par Por su cuenta, en el par de impágenes de la Figura \ref{fig:NSURF} se muestran emparejamientos nulos con el algoritmo, estos se deben a la orientación del señalamiento en la imagen.
\begin{figure}[htbp!]
	\centering
	\begin{subfigure}[htbp!]{0.4\textwidth}
		\includegraphics[width=\textwidth]{./Figuras/NSURF1}
	\end{subfigure}
	\begin{subfigure}[htbp!]{0.4\textwidth}
		\includegraphics[width=\textwidth]{./Figuras/NSURF2}
	\end{subfigure}
	\caption{Emparejamientos nulos con el algoritmo SURF.}
	\label{fig:NSURF}
\end{figure}
\par Por último, para seguir la navegación del vehículo durante todo su recorrido se implementó la odometría visual basado en descriptores FAST. Las funciones concernientes a la odometría visual están presentes en la clase Navegación. Como resultado de esta implementación, en la Figura \ref{fig:odom} se muestra la salida del programa para dos pruebas en las que el programa trabajó {\it fuera de línea}; el programa devuelve una ventana de 300x300 pixeles que representa un escenario de 3x3 m. En la ventana se observa a lo largo del tiempo cómo se dibujan nuevos puntos azules, indicando de manera aproximada el camino recorrido por el vehículo desde su configuración inicial a su configuración final.
\begin{figure}[htbp!]
	\centering
	\begin{subfigure}[htbp!]{0.4\textwidth}
		\includegraphics[width=\textwidth]{./Figuras/Odo1}
	\end{subfigure}
	\begin{subfigure}[htbp!]{0.4\textwidth}
		\includegraphics[width=\textwidth]{./Figuras/Odo2}
	\end{subfigure}
	\caption{Seguimiento de la trayectoria del vehículo a través de odometría visual.}
	\label{fig:odom}
\end{figure}
\par Estas imágenes finales correspondientes a la odometría, por su caracter de trabajo extra, no se encuentran documentadas en su totalidad, y en la Sección de trabajo futuro se establece que se continuará trabajando en este elemento.
	\chapter{Conclusiones y trabajo futuro}
	\label{cap:con}
	Desde el Protocolo de Investigación se plantearon objetivos cuyo cumplimiento se puede clasificar en cuatro vertientes: cumplido satisfactoriamente, cumplido insatisfactoriamente, cumplido parcialmente y no cumplido. Se puede concluir, después de haber analizando cada uno de los objetivos por separado, lo siguiente:
\par El objetivo general dice: {\it Implementar algoritmos de visión artificial para la identificación de caracterticas del entorno, que serán usadas como realimentación visual en un esquema de control para la navegación de un vehículo autónomo a escala en un escenario urbano simulado}. Al término de este trabajo terminal se puede concluir que el objetivo ha sido \underline{cumplido satisfacoriamente}, y de él emanan las siguientes conclusiones con respecto a los objetivos específicos.
\begin{itemize}
	\item {\it Escoger una plataforma robótica de propulsión eléctrica para validar la navegación autónoma del proyecto.} Este objetivo se \underline{cumplió satisfactoriamente} desde TT1, y se puede realizar esta afirmación ya que para los requerimientos de este proyecto no es necesario contar con un chasís de alta resistencia en diferentes condiciones, tampoco es necesario que pueda desplazarse a altas velocidades, o posea cualidades que le permitan desplazarse en terrenos accidentados. Para este punto sólo es necesario que el vehículo seleccionado sea un vehículo a escala 1:10, por lo que el haber seleccionado una plataforma económica no repercutió en el desarrollo.
	\item {\it Seleccionar los algoritmos de visión artificial adecuados para el proyecto.} Con relación a este objetivo, se basó el desarrollo en lo encontrado en el Estado del Arte sin ir más allá, por lo que se \underline{cumplió a secas} el objetivo.
	\item {\it Realizar el análisis cinemático del robot móvil para la obtención de parámetros que luego se emplearán en el control del vehículo.} Como se vio en la Sección \ref{sec:mc}, se esudió el modelo cinemático del vehículo para tomar en cuenta distintos parámetros importantes a tomar en cuenta para el controlador; después vienen las simulaciones donde se comprueba el comportamiento del vehículo al pasar de una configuración inicial a una configuración final con ayuda de un controlador; es por eso que se considera que este objetivo se ha {\it cumplido satisfactoriamente}.
	\item {\it Diseñar y programar los algoritmos de control visual para su implementación en el robot móvil.}
	\item {\it Usar técnicas de control para regular la velocidad y dirección del vehículo autónomo a escala.} El producto más importante de estos dos objetivos es el controlador lateral del vehículo que le permite mantenerse en su carril mientras navega en la pista de pruebas,. El objetivo, por tanto, se encuentra \underline{cumplido satisfactoriamente}.
	\item {\it Emplear la realimentación visual para alcanzar un nivel SAE 3 de automatización \cite{saeTaxonomyDefinitiosTerms2014} en el vehículo a escala.} Con los resultados obtenidos para este proyecto, este objetivo se encuentra \underline{no cumplido}, pues el vehículo no incorpora forma alguna para identificación de obstáculos y no se logra la navegación continua del mismo.
	\item {\it Construir un escenario de pruebas para la validación de la plataforma a través de: a) Identificar y permanecer en el carril por el que se desplaza; b) Reconocer las eñales de Alto, vuelta a la Izquierda y vuelta a la Derecha en el escenario de pruebas; c) Examinar el color de la luz en un semáforo y reaccionar de manera correcta ante ello.} En la Sección \ref{sec:platpru} se muestra el proceso de construcción de la pista de pruebas, a la vez que los objetivos específicos cuarto y quinto se relacionan con el inciso a) de este objetivo; los incisos b) y c), como se encuentra en el Capítulo \ref{cap:resul}, la localización de los señalamientos sólo es apta para funcionar en condiciones específicas y son susceptibles de cambios ambientales, por lo que se considera que el objetivo ha sido \underline{cumplido de manera instisfactoria.}
\end{itemize}
En general, se logró cumplir con los objetivos presentados como retos para la plataforma robótica, de manera que como posible desarrollo posterior se presena a continuación el trabajo futuro.
\section{Trabajo futuro}
\label{sec:trafut}
Las posibles futuras mejoras y adaptaciones a la plataforma robótica se plantean las siguientes:
\begin{itemize}
	\item Incorporar nuevos sensores para el vehículo que permitan otorgarle una clasificación de autonomía mayor según \cite{saeTaxonomyDefinitiosTerms2014}.
	\item Explorar diferentes sistemas de locomoción para las mismas acciones realizadas en este trabajo.
	\item Buscar nuevos algoritmos de visión artificial que vayan más allá de los tomados del estado del arte.
	\item Poner en funcionamiento un facsímil de este proyecto sobre ROS.
	\item Mejorar el controlador implementado en este trabajo.
	\item Implementar en el robot controladores difusos y de redes neuronales.
	\item Mejorar la pista de pruebas para que permita al vehículo planear trayectorias.
	\item Mejorar la pista de pruebas para que permita maniobras grupales (yendo más allá de las individuales tomadas exploradas en este documento).
	\item Escalar el vehículo en cuanto a comunicación para que pueda permitir las maniobras grupales con varios vehículos.
\end{itemize}
	\addcontentsline{toc}{chapter}{Bibliografía}
	\bibliographystyle{IEEEtran}
	\bibliography{Biblio}
	\label{chap:Biblio}
	\appendix
	\addcontentsline{toc}{part}{Anexos}
	\chapter{Hojas de datos y manuales}
	\label{chap:datasheet}
	\vspace*{-0.5cm}
	\section{Manual de usuario del ESC}
	\label{sec:ESC}
	\vspace*{-1cm}
	\begin{figure}[htbp!]
		\centering
		\includegraphics[width=0.8\textwidth, angle=180]{./Figuras/001}
	\end{figure}
	\begin{figure}[htbp!]
		\centering
		\includegraphics[width=0.9\textwidth, angle=180]{./Figuras/002}
	\end{figure}
	\newpage 
	\section{Manual de usuario del vehículo a escala}
	\label{sec:vrx}
	\vspace*{-0.3cm}
	\begin{figure}[htbp!]
		\centering
		\includegraphics[width=0.8\textwidth, angle=180]{./Figuras/003}
	\end{figure}
	\newpage
	\section{GPIO de la computadora Rock64}
	\label{sec:pi}
	\begin{figure}[htbp!]
		\centering
		\includegraphics[width=0.8\textwidth]{./Figuras/Pi-2}
	\end{figure}
	\begin{figure}[htbp!]
		\centering
		\includegraphics[width=0.8\textwidth]{./Figuras/Pi-5}
	\end{figure}
	\chapter{Planos de construcción}
	\label{chap:plan}
	\vspace*{-0.5cm}
	\section{Chasís superior}
	\label{sec:chass}
	\vspace*{-0.5cm}
	\begin{figure}[htbp!]
		\centering
		\includegraphics[width=1.1\textwidth, angle=270]{./Figuras/Placa.PDF}
	\end{figure}
	\chapter{Puesta a punto de la computadora Rock64}
	\label{cap:pap}
	Para este proceso se requerirá de una PC y la Rock64; este apéndice desarrolla los pasos necesarios para transferir la imagen del sistema operativo desde la PC a la Rock64 y posteriormente instalar en ésta los paquetes necesarios para el desarrollo del proyecto.
	\section{Instalación del sistema operativo en la memoria microSD}
	La compañía Pine64 provee una herramienta para gabar el sistema operativo elegido en la memoria microSD. Una vez instalado el programa en la PC desde el repositorio\footnote{https://github.com/pine64dev/PINE64-Installer}, en el programa se elige para cuál producto se hará el proceso y después será necesario elegir el sistema operativo, resultando la pantalla siguiente. 
	\begin{figure}[htbp!]
		\centering
		\includegraphics[width=0.8\textwidth]{./Figuras/Etcher}
		\caption{Instalación del sistema operativo en la memoria microSD}
		\label{fig:etcher}
	\end{figure}
	\par Una vez realizado esto, se elige la ruta de la memoria microSD a grabar y se presiona el botón Flash! Para este proyecto la distribución elegida es ARMbian Bionic Desktop y este proceso toma algunos minutos en completarse, tras lo cual se puede proceder a conectar la memoria microSD en la Rock64, con lo que en el primer arranque se obtendrá la pantalla de la figura \ref{fig:ARMbian}.
	\begin{figure}[htbp!]
		\centering
		\includegraphics[width=0.8\textwidth]{./Figuras/Boot}
		\caption{Pantalla del ARMbian en el primer arranque}
		\label{fig:ARMbian}
	\end{figure}
	\par Tras ingresar credenciales se inicia el sistema operativo en modo gráfico (porque esta opción fue la seleccionada), y ahora se pueden realizar los siguientes pasos para la puesta a punto de la tarjeta.
	\section{Instalación de dependencias}
	La siguiente serie de instrucciones se realiza en la terminal de la computadora. Se puede ingresar a ella navegando con ayuda del ratón y también es posible tecleando Ctrl+Alt+T. Una vez abierta teclear los siguiente comandos en orden.
	\lstset{language=bash}	
	\begin{lstlisting}[frame=single]
	sudo apt-get purge galculator* leafpad* geany* mirage* hexchat* pulseaudio* atril* orage* 
	sudo apt-get update
	sudo apt-get upgrade
	sudo apt-get dist-upgrade
	sudo apt-get install emacs
	sudo apt-get install tightvncserver install openjdk-8-jdk
	sudo apt-get install 12c-tools libc6-dev man manpages manpages-dev libi2c-dev python-smbus libasound2-dev libasound2-doc sqlite2 libsqlite3-dev sqlite3-doc zlib1g-dev libX11-dev libxcb-doc less gnuplot ntp bc
	sudo apt-get install git libssl-dev libusb-1.0.0-dev libudev-dev pkg-config libgtk-3-dev libglfw3-dev freeglut3-dev make cmake cmake-gui build-essential unzip libxmu-dev libxi-dev
	sudo apt-get install libjpeg-dev libpng-dev libtiff-dev libavcodec-dev libavformat-dev libswscale-dev libv4l-dev libxvidcore-dev libx264-dev
	sudo apt-get install libgtk-3-dev
	sudo apt-get install libatlas-base-dev gfortran
	sudo apt-get install python3-dev
	wget https://bootstrap.pypa.io/get-pip.py
	sudo python3 get-pip.py
	sudo pip install virtualenv virtualenvwrapper
	sudo rm -rf ~/get-pip.py ~/.cache/pip
	sudo apt-get install libopenni2-dev
	\end{lstlisting}
	Después de esto, se ingresa al archivo /.bashrc por medio del editor de texto que se quiera y al final del archivo se escribe el siguiente código:
	\begin{lstlisting}[frame=single]
	# virtualenv and virtualenvwrapper
	export WORKON_HOME=$HOME/.virtualenvs
	export VIRTUALENVWRAPPER_PYTHON=/usr/bin/python3
	source /usr/local/bin/virtualenvwrapper.sh
	\end{lstlisting}
	Se ejecuta el comando y se configuran nueva dependencias al usar los siguientes comandos
	\begin{lstlisting}[frame=single]
	source ~/.bashrc
	mkvirtualenv cv -p python3
	workon cv
	pip install numpy
	\end{lstlisting}
	\section{Instalación de LibRealSense}
	Si se ha descargado y descomprimido previamente la librería, entonces se ejecutan las siguientes órdenes:
	\begin{lstlisting}[frame=single]
	cd ~/librealsense
	sudo ./scripts/setup_udev_rules.sh
	sudo ./scripts/patch-realsense-ubuntu-lts.sh
	mkdir build && cd build
	cmake -D BUILD_OPENNI2_BINDINGS=true -DOPENNI2_DIR=/usr/include/openni2 -DREALSENSE2_DIR=/usr/local/include/librealsense2 -DFORCE_LIBUVC=true ..
	sudo make uninstall && make clean && make -j4 && sudo make install
	sudo ldconfig
	\end{lstlisting}
	Para esta librería y también para OpenCV, la configuración de la orden cmake se puede realizar también desde la cmake-gui. De la misma manera, el proceo de compilación para ambas librerías es tardado, alrededor de tres horas para ambas.
	\section{Instalación de OpenCV}
	\begin{lstlisting}[frame=single]
	cd ~/opencv
	mkdir build && cd build
	cmake -DCMAKE_BUILD_TYPE=RELEASE -DBUILD_EXAMPLES=ON -DCMAKE_INSTALL_PREFIX=/usr/local -DWITH_OPENNI2=true -DWITH_LIBREALSENSE=true ..
	make -j4
	sudo make install
	sudo ldconfig
	\end{lstlisting}
	\chapter{Códigos de programación}
	\label{cap:cod}
	\section{Simulaciones}
	\label{sec:csim}
	\lstinputlisting[language=Python, caption={Programa usado para las simulaciones (Python)}]{./Programas/Simulaciones/PID.py}
	\newpage
	\section{Programa del microcontrolador}
	\label{sec:micro}
	\lstinputlisting[language=C++, caption={Programa grabado en el microcontrolador (C)}]{./Programas/AVR/Servo.c}
	\newpage
	\section{Clase Cámara}
	\label{sec:cam}
	\lstinputlisting[language=C++, caption={Clase Cámara, archivo de cabecera (C++)}]{./Programas/Proyecto/include/cam.hpp}
	\lstinputlisting[language=C++, caption={Clase Cámara, archivo de funciones (C++)}]{./Programas/Proyecto/src/cam.cpp}
	\section{Clase Control}
	\label{sec:cont}
	\lstinputlisting[language=C++, caption={Clase Control, archivo de cabecera (C++)}]{./Programas/Proyecto/include/control.hpp}
	\lstinputlisting[language=C++, caption={Clase Control, archivo de funciones (C++)}]{./Programas/Proyecto/src/control.cpp}
	\section{Clase Carro}
	\label{sec:car}
	\lstinputlisting[language=C++, caption={Clase Carro, archivo de cabecera (C++)}]{./Programas/Proyecto/include/carro.hpp}
	\lstinputlisting[language=C++, caption={Clase Carro, archivo de funciones (C++)}]{./Programas/Proyecto/src/carro.cpp}
	\section{Programa principal}
	\label{sec:main}
	\lstinputlisting[language=C++, caption={Programa principal (C++)}]{./Programas/Proyecto/src/main.cpp}
\end{document}